\documentclass[11pt,a4paper]{article}
\usepackage[top=2.5cm,bottom=2.8cm,left=2.5cm,right=2.5cm]{geometry}
\usepackage{amsmath,amssymb,amsthm}
\usepackage{hyperref,xcolor,enumitem,booktabs,array}
\usepackage{graphicx,subcaption}

\hypersetup{colorlinks=true,linkcolor={blue!60!black},
  citecolor={blue!60!black},urlcolor={blue!60!black}}

\newtheorem{theorem}{Theorem}[section]
\newtheorem{lemma}[theorem]{Lemma}
\newtheorem{proposition}[theorem]{Proposition}
\newtheorem{corollary}[theorem]{Corollary}
\theoremstyle{definition}\newtheorem{definition}[theorem]{Definition}
\theoremstyle{remark}
\newtheorem{observation}[theorem]{Observation}
\newtheorem{remark}[theorem]{Remark}
\newtheorem{conjecture}[theorem]{Conjecture}

\newcommand{\vp}{\varphi}
\newcommand{\lam}{\lambda}
\newcommand{\Ja}{J_{2a}}
\newcommand{\Jfour}{J_4}
\newcommand{\TCMB}{T_{\mathrm{CMB}}}
\newcommand{\Lcond}{\Lambda_{\mathrm{cond}}}
\newcommand{\rCMB}{\rho_{\mathrm{CMB}}}
\newcommand{\SU}{\mathrm{SU}}
\newcommand{\Efb}{E_{\mathrm{fb}}}
\newcommand{\Kfb}{K_{\mathrm{fb}}}
\newcommand{\twoI}{2I}
\newcommand{\twoT}{2T}
\newcommand{\QH}{Q_H}
\newcommand{\vEW}{v_{\mathrm{EW}}}
\newcommand{\Ztwo}{\mathbb{Z}_2}
\newcommand{\Zthree}{\mathbb{Z}_3}
\newcommand{\dimq}[1]{d_{#1}}
\newcommand{\Tmat}{T}
\newcommand{\Jones}{V}
\newcommand{\lk}{\mathrm{lk}}
\newcommand{\phistar}{\phi^*}

\title{\textbf{Preimage Topology and Confinement:\\
The $Q_H=3$ Sector of the Density-Feedback Hopfion}
\\[0.6em]
{\large Companion Paper XVIII to the Density-Feedback
Faddeev--Niemi Hopfion Series}
\\[0.6em]
\large DRAFT --- June 2026}
\author{Frederick Manfredi}
\date{}

\begin{document}
\maketitle

\begin{abstract}
Papers~XV and~XVI establish the $\QH=3$ (trefoil $T_{2,3}$, binary
tetrahedral group $\twoT$) sector as the baryon sector and identify
confinement with the simple-current condition $\dimq{3/2}=1$ at
WZW level $k=3$.  Paper~XVII further resolves the Dirac/Majorana
question for the $\QH=1$ neutrino sector via the Wess--Zumino descent.
The present paper develops the \emph{dynamical} content of $\QH=3$
confinement by analysing the Hopf preimage topology of the torus
knot/link tower and deriving three results.
(1)~\textbf{Confinement from Jones polynomial arithmetic: fusion, not sector, splits rational from irrational.}
The satellite Jones invariant of the $\QH=3$ trefoil (computed via
the $(2,3)$-cable/companion construction of Section~\ref{P18:sec:satellite})
satisfies $|\Jones_{\rm sat}(q_5)|=1/\vp$ at the WZW root $q_5=e^{2\pi i/5}$,
and the Hopf link's invariant satisfies $|\Jones_{T(2,2)}(q_5)|=1/\vp$
also: the two sectors carry the \emph{same} modulus, related by the
single phase $\Jones_{\rm sat}(q_5)=q_5^{-1}\cdot\Jones_{T(2,2)}(q_5)$,
not by an inversion of the golden ratio.
A satellite twist family (Proposition~\ref{P18:prop:satellite_twist_family})
reveals where the rational/irrational split actually lies: every
\emph{fused} (genuine torus knot, odd half-twist count $n$)
configuration has an irrational $\vp$-related modulus, while every
\emph{unfused} (cables not yet merged, even $n$) configuration has a
rational invariant ($V=1$ at $n=4$, $V=-1$ at $n=6$) --- the split
tracks fusion of the cable strands, not which $\QH$ sector is
involved. Whether the physical formation path passes through the
rational unfused configuration at $n=4$, and whether the family's
period ($10$) coincides with the quantum-group order
$Q_{\mathrm{gr}}=2(k+2)=10$, are noted as open conjectures
(Section~\ref{P18:sec:jones}), not established results.
(2)~\textbf{Forced pion emission.}  The reaction
$\QH=2+\QH=2\to\QH=3+\QH=1$ is topologically forced by charge
conservation ($2+2=3+1$) and the component structure of the
satellite construction (Proposition~\ref{P18:prop:satellite}):
of the four initial circles, three fuse into the trefoil and
one cannot be absorbed, producing the $\QH=1$ spectator.
Pion emission in baryon formation is therefore a topological
necessity, not a dynamical accident.
(3)~\textbf{Process-dependent $\phistar$ prediction.}  The three
crossing vertices of $T_{2,3}$ are 120° apart in azimuth and their
tangent vectors lie exactly in the $z=0$ plane.  When a baryon
($\QH=3$) fragments to produce $\rho^+\rho^-$ ($\QH=2$ pairs), the
relative decay-plane angle $\phistar$ is set by this crossing geometry,
predicting a distribution $W(\phistar)\propto 1+a_3\cos(3\phistar)$
peaked at $\phistar=2\pi/3$, with $a_1=a_2=0$ enforced by $\Zthree$
symmetry.  Photon-initiated $\rho^+\rho^-$ production (no $\QH=3$
parent) carries no such geometric constraint, giving the standard
helicity-governed $a_2\cos(2\phistar)$ form.  The process-dependence
is a parameter-free, falsifiable prediction.
Two speculative correspondences --- a six-quark counting from
$\binom{3}{1}+\binom{3}{2}=6$ midpoint-breaking combinations, and an
isospin assignment from the $\Ztwo$ writhe asymmetry of $T_{2,3}$ ---
are recorded as conjectures.
\end{abstract}

\tableofcontents
\bigskip\newpage

%══════════════════════════════════════════════════════════════════════════════
\section{Introduction}
\label{P18:sec:intro}
%══════════════════════════════════════════════════════════════════════════════

The density-feedback Faddeev--Niemi (FN) Hopfion framework
assigns each topological sector $\QH=N$ its own identity through a
McKay-correspondence group, a WZW conformal field theory, and a
knot or link geometry.  Through Paper~XVII~\cite{Paper17}, the three
matter sectors are settled at the qualitative level:
\begin{center}
\renewcommand{\arraystretch}{1.35}
\begin{tabular}{lllllll}
\toprule
$\QH$ & Particle & Group & McKay & WZW & $Q_{\mathrm{gr}}$ & $\dimq{}$ \\
\midrule
$1$ & Neutrino & $\Ztwo$ & $\widehat{A}_1$ & $\SU(2)_1$ & $6$ & $\vp$ \\
$2$ & Charged lepton & $\twoI$ & $\widehat{E}_8$ & $\SU(2)_3$ & $10$ & $\vp$ \\
$3$ & Baryon & $\twoT$ & $\widehat{E}_6$ & $(E_6)_1$ & $3$ & $1$ \\
\bottomrule
\end{tabular}
\end{center}
What remains missing is the \emph{dynamical} content of confinement
in the $\QH=3$ sector: why does the simple-current condition
$\dimq{3/2}=1$ manifest as string-breaking, what topological reaction
rules govern the production of $\QH=2$ and $\QH=1$ objects from a
$\QH=3$ parent, and what experimental signatures follow?

This paper answers those questions in three stages.

\paragraph{Stage 1 (Sections~\ref{P18:sec:jones}--\ref{P18:sec:preimage}).}
We establish the Jones polynomial as the correct knot-invariant
framework for the sector hierarchy.  The satellite Jones invariant
of the $T(2,n)$ tower at the WZW root $q_5=e^{2\pi i/5}$ exhibits
a precise structure: every fused (genuine torus knot) member carries
an irrational modulus, while every unfused (cables not yet merged)
member carries a rational value.  Both $|\Jones_{T(2,2)}(q_5)|$
($\QH=2$, Hopf link) and $|\Jones_{\rm sat}(T(2,3))(q_5)|$
($\QH=3$, trefoil satellite) equal $1/\vp$, related by the single
phase $q_5^{-1}$; the rational/irrational distinction tracks fusion
of the cable strands, not the $\QH=2\to\QH=3$ transition itself.
The Verlinde quantum dimension condition $\dimq{3/2}=1$
(Paper~XVII~\cite{Paper17}, Proposition~7.2 %\ref{P17:prop:dimq_confinement}
) remains the independent, established confinement criterion; how it
relates to the fused/unfused split of the satellite family is
discussed but not resolved.

\paragraph{Stage 2 (Sections~\ref{P18:sec:preimage}--\ref{P18:sec:pion}).}
We analyse the Hopf preimage structure of each sector.  For $\QH=2$,
the preimage of a regular value on $S^2$ under the Hopf map is the
two-component Hopf link $T(2,2)$.  For $\QH=3$, the correct
preimage configuration is the trefoil $T(2,3)$ (a single knotted
circle), reached from $T(2,2)$ by the satellite/cable construction.
This determines the topological reaction rules and, in particular,
forces pion emission ($\QH=1$) as a necessary companion of baryon
($\QH=3$) formation from two $\QH=2$ objects.

\paragraph{Stage 3 (Sections~\ref{P18:sec:phi_star}--\ref{P18:sec:arcseg}).}
We derive the process-dependent $\phistar$ prediction from the exact
geometry of $T_{2,3}$'s crossing vertices, connect it to the
arc-segment energy partition of Paper~XVI~\cite{Paper16}, and
identify the experimental test that distinguishes baryon-initiated
from photon-initiated $\rho^+\rho^-$ production.

Sections~\ref{P18:sec:quark_count}--\ref{P18:sec:isospin} collect speculative
correspondences, clearly flagged as conjectures, that connect the
12-site fraying pattern to quark flavour counting and the $\Ztwo$
writhe asymmetry to isospin.

%══════════════════════════════════════════════════════════════════════════════
\section{The Jones Polynomial at the WZW Root of Unity}
\label{P18:sec:jones}
%══════════════════════════════════════════════════════════════════════════════

\subsection{The WZW root of unity}
\label{P18:sec:root_of_unity}

The $\SU(2)_k$ WZW model at level $k$ has modular data governed by
the $q$-deformation parameter $q = e^{i\pi/(k+2)}$.  At $k=3$,
\begin{equation}
  q = e^{i\pi/5}, \qquad q^2 = e^{2\pi i/5}.
  \label{P18:eq:q_wzw}
\end{equation}
This is a primitive 10th root of unity (equivalently a primitive
5th root of $q^2$).  The quantum integers $[n]_q =
(q^n-q^{-n})/(q-q^{-1})$ at this value satisfy $[5]_q = 0$
(the truncation that gives the finite WZW representation ring),
and the quantum dimension of the spin-$j$ representation is
\begin{equation}
  \dimq{j} = [2j+1]_q = \frac{\sin((2j+1)\pi/5)}{\sin(\pi/5)},
  \label{P18:eq:dimq}
\end{equation}
recovering Table~1 of Paper~XVII~\cite{Paper17} at $k=3$:
$\dimq{1/2} = \vp$, $\dimq{1} = \vp$, $\dimq{3/2} = 1$,
$\dimq{2} = 0$ (the null representation, boundary of the representation
ring).

\subsection{Jones polynomials of the torus link/knot tower}

The Jones polynomial $\Jones_L(q)$ of a link $L$, evaluated at
$q = e^{2\pi i/5}$ (the squared WZW root), encodes exactly the
quantum-dimension data of the WZW representation theory.

\begin{theorem}[Satellite Jones invariant of the trefoil]
\label{P18:thm:trefoil_satellite_jones}
Let $q_5=e^{2\pi i/5}$.  The trefoil $T(2,3)$ is not a free-standing
knot in this framework: by Proposition~\ref{P18:prop:satellite} it is the
$(2,3)$-cable of one component of the $\QH=2$ Hopf link, with the
other (companion) component still present and linked.  Computing the
Jones invariant of this configuration correctly --- via the
Kauffman-bracket/Temperley--Lieb cabling formula applied to the
diagram with the companion included, rather than to the trefoil taken
as an abstract free-standing knot --- gives
\begin{equation}
  \Jones_{\mathrm{sat}}(q_5) \;=\; \frac{q_5^3}{\vp},
  \qquad
  |\Jones_{\mathrm{sat}}(q_5)| \;=\; \frac{1}{\vp}.
  \label{P18:eq:trefoil_satellite_value}
\end{equation}
This is irrational: the golden ratio is present in the trefoil's
satellite invariant at $q_5$, not absent from it.
\end{theorem}

\begin{proof}
See \texttt{papers/src\_paper18/jones\_T23\_satellite\_check.py}. The
calculation engine (a first-principles Kauffman-bracket/Temperley-Lieb
calculator, \texttt{kauffman\_bracket\_engine.py} in the same
directory) is validated by exact reproduction of the standard
tabulated Jones polynomials of the free trefoil
($V(t)=-t^4+t^3+t$) and the free Hopf link
($V(t)=-t^{1/2}-t^{5/2}$) before being applied to the satellite
diagram (cable strands twisted $\sigma_1^3$, companion linked once
with each cable strand, matching the physical construction of
Proposition~\ref{P18:prop:satellite}).  The resulting link has two
components (the fused cable pair and the separate companion),
verified by an independent permutation/component-count check in the
same script.
\qed
\end{proof}

\begin{remark}[The free-standing evaluation, for comparison]
\label{P18:rem:trefoil_satellite_correction}
Evaluated as an abstract, free-standing knot (i.e.\ ignoring the
satellite embedding of Proposition~\ref{P18:prop:satellite}), the
trefoil's Jones polynomial at $q_5$ is $-q_5^{-1}+q_5^3+q_5$, of
modulus $\sqrt{4-\vp}$ --- a different irrational value from the
satellite invariant of Theorem~\ref{P18:thm:trefoil_satellite_jones}.
Neither is an integer, and the golden ratio is present in both.
Whether and how the trefoil's departure from the integer value $1$
(the Jones invariant of the $\QH=2$ Hopf link's own unknotted
companion, prior to cabling) relates to the transience of the
$\QH=3$ sector is not yet established and is left open.
\end{remark}

\begin{proposition}[The trefoil within the satellite twist family]
\label{P18:prop:satellite_twist_family}
Consider the family of satellite invariants obtained by varying the
number of half-twists $n$ between the two cable strands, with the
companion linked once to each cable strand throughout as in
Theorem~\ref{P18:thm:trefoil_satellite_jones} (the case $n=3$).  The
modulus $|\Jones_{\mathrm{sat}}^{(n)}(q_5)|$ is periodic in $n$ with
period $10$, and satisfies the mirror symmetry
$|\Jones_{\mathrm{sat}}^{(n)}(q_5)|=|\Jones_{\mathrm{sat}}^{(10-n)}(q_5)|$.
At the four smallest cases:
\begin{align}
  n=3\ (\text{trefoil }T(2,3),\ \QH=3,\ \text{fused}):
    &\quad |\Jones_{\mathrm{sat}}(q_5)| = 1/\vp,
    \label{P18:eq:sat_n3}\\
  n=4\ (\text{unfused, three components}):
    &\quad \Jones_{\mathrm{sat}}(q_5) = 1,
    \label{P18:eq:sat_n4}\\
  n=5\ (\text{cinquefoil } T(2,5),\ \text{fused}):
    &\quad |\Jones_{\mathrm{sat}}(q_5)| = \vp,
    \label{P18:eq:sat_n5}\\
  n=6\ (\text{unfused, three components}):
    &\quad \Jones_{\mathrm{sat}}(q_5) = -1.
    \label{P18:eq:sat_n6}
\end{align}
Every fused configuration computed (odd $n$, a genuine torus knot)
gives an irrational, golden-ratio-related modulus; every unfused
configuration computed (even $n$, two separate, un-merged cable
strands) gives a rational value.
\end{proposition}

\begin{proof}
See \texttt{papers/src\_paper18/jones\_satellite\_twist\_family.py},
which extends the calculation of
Theorem~\ref{P18:thm:trefoil_satellite_jones} to general twist number $n$
and verifies the stated period, symmetry, and exact values
symbolically.
\end{proof}

\begin{remark}[Where the integer value $-1$ actually occurs]
\label{P18:rem:minus_one_at_n6}
The value $\Jones(q_5)=-1$ originally claimed in this section for the
trefoil does occur in this family --- exactly, at $n=6$, an unfused
three-component configuration --- rather than at $n=3$.
\end{remark}

\begin{remark}[The confinement criterion: fusion, not sector, and a phase relation between $\QH=2$ and $\QH=3$]
\label{P18:rem:richer_criterion}
The satellite twist family of Proposition~\ref{P18:prop:satellite_twist_family},
together with a direct computation of the (uncabled) Hopf link's own
Jones invariant, gives a confinement criterion that is different
from --- and more subtle than --- the original integer/irrational
framing.  We record three pieces of structure and one open question.

\textbf{(A) Confinement marks fusion, not sector.}
The rationality-class split tracks whether the cable strands
have \emph{fused} into a genuine torus knot, not which topological
sector is involved:
\begin{center}
\renewcommand{\arraystretch}{1.3}
\begin{tabular}{ccccl}
\toprule
$n$ & Components & $|\Jones_{\rm sat}(q_5)|$ & Rational? & Notes \\
\midrule
3 & 2 (fused, $T(2,3)$) & $1/\vp$ & no & $\QH=3$ trefoil \\
4 & 3 (unfused)          & $1$ & \textbf{yes} & \\
5 & 2 (fused, $T(2,5)$) & $\vp$ & no & cinquefoil \\
6 & 3 (unfused)          & $1$ & \textbf{yes} & [$V=-1$, so rational] \\
\bottomrule
\end{tabular}
\end{center}
Every fused (odd $n$) state computed is irrational; every unfused
(even $n$) state computed is rational.  Whether the physical
formation path from $\QH=2$ to $\QH=3$ actually passes through the
$n=4$ configuration, and whether its rationality is what produces the
flat-bottomed energy landscape of Paper~XVI~\cite{Paper16}
(Remark~\ref{P18:rem:saddle_difficulty}), is an open conjecture, not
established: nothing in this section identifies the twist parameter
$n$ with a physically continuous coordinate of the actual
condensate dynamics.

\textbf{(B) The Hopf link and the trefoil satellite carry the same
modulus, related by a phase.}  Direct computation of the (uncabled)
Hopf link's own Jones polynomial (not the disjoint-union factor, a
different quantity discussed below) gives
\begin{equation}
  |\Jones_{T(2,2)}(q_5)| \;=\; \frac{1}{\vp}
  \;=\; |\Jones_{\rm sat}(T(2,3))(q_5)|,
  \qquad
  \Jones_{\rm sat}(q_5) \;=\; q_5^{-1}\cdot\Jones_{T(2,2)}(q_5).
  \label{P18:eq:phase_relation}
\end{equation}
The two sectors carry the \emph{same} modulus; they are distinguished
by a single unit of the phase $q_5$, not by an inversion of the
golden ratio.  Whether this phase relation has independent physical
content, or how (if at all) it connects to the simple-current
condition $\dimq{3/2}=1$ (an established, independent confinement
criterion from Paper~XVII~\cite{Paper17}, Proposition~7.2), is not
established here.
(The disjoint-union factor $-(q_5^{1/2}+q_5^{-1/2})=-\vp$, used
elsewhere in this paper for multi-component reactions,
Theorem~\ref{P18:thm:jones_rule}, is a different quantity from
$\Jones_{T(2,2)}(q_5)$ and should not be conflated with it.)

\textbf{(C) The tower structure is $\{1/\vp, 1, \vp, -1, \ldots\}$.}
The fused states (odd $n$) have moduli alternating between $1/\vp$
and $\vp$ under the mirror symmetry
$|\Jones_{\rm sat}^{(n)}|=|\Jones_{\rm sat}^{(10-n)}|$;
the unfused states (even $n$) alternate between $V=1$ and $V=-1$.
This period is $10$; whether this is the same $10$ as the
quantum-group order $Q_{\rm gr}=2(k+2)=10$ of the framework
(Paper~I~\cite{Paper1}), or a coincidence of this specific tangle's
eigenvalue orders, is an open question and is not established.

\textbf{(D) Where $-1$ lives: the unfused sector at $n=6$.}
The value $V=-1$ that was originally claimed for the trefoil ($n=3$)
does occur in the family, but at $n=6$ (an unfused, three-component
configuration, the next rational step beyond the cinquefoil).
The classical identity $V_{T(2,3)}(i)=-1$ holds at the 4th root of
unity $t=i=e^{2\pi i/4}$ (the $k=2$ WZW level, the Arf invariant
identity $V_K(i)=(-1)^{{\rm Arf}(K)}$), which is unrelated to
the $k=3$ WZW structure used throughout this framework.
The precise placement of $-1$ at $n=6$ (unfused, rational) rather
than $n=3$ (fused, irrational) is itself consistent with the
fused/unfused rule of point~(A).
\end{remark}

%══════════════════════════════════════════════════════════════════════════════
\section{Hopf Preimage Structure and the Torus Knot Tower}
\label{P18:sec:preimage}
%══════════════════════════════════════════════════════════════════════════════

\subsection{Preimage circles and the Hopf fibration}
\label{P18:sec:preimage_hopf}

For a smooth map $\mathbf{n}:\mathbb{R}^3\to S^2$ in the homotopy
class $\QH=N$ (Hopf degree $N$), the preimage $\mathbf{n}^{-1}(p)$
of a regular value $p\in S^2$ is a disjoint union of smooth oriented
circles $\{\gamma_a\}$ in $\mathbb{R}^3$.  The Hopf invariant is
recoverable from the total pairwise linking number:
\begin{equation}
  \QH \;=\; \sum_{a<b} \lk(\gamma_a,\gamma_b),
  \label{P18:eq:hopf_link}
\end{equation}
where $\lk$ denotes the Gauss linking number.

\begin{definition}[Torus knot/link tower]
\label{P18:def:torus_tower}
The \emph{torus knot/link $T(2,n)$} is defined as the closure of
the two-strand braid $\sigma_1^n\in B_2$, embedded on the standard
torus $\mathbb{T}^2\hookrightarrow S^3$ as the curve that winds $2$
times around the longitude and $n$ times around the meridian.
For $n$ odd, $T(2,n)$ is a \emph{knot} (one component).
For $n$ even, $T(2,n)$ is a \emph{link} (two components, linking
number $n/2$).
\end{definition}

\begin{proposition}[Preimage topology of the $T(2,n)$ tower]
\label{P18:prop:preimage_tower}
In the density-feedback FN framework, the minimal-action condensate
configuration in sector $\QH=N$ has the following preimage topology:
\begin{enumerate}[noitemsep,label=(\roman*)]
\item $\QH=1$: preimage is a single unknotted circle $T(2,1)$.
\item $\QH=2$: preimage is the two-component Hopf link $T(2,2)$,
  with $\lk(\gamma_1,\gamma_2)=1$ and the single
  pair gives $\QH=2\cdot\lk=2$.
\item $\QH=3$: preimage is the trefoil knot $T(2,3)$, a single
  knotted circle.  The three self-crossings of $T(2,3)$ in its
  standard diagram correspond structurally to the three quark
  crossing vertices of Paper~XV--XVI.
\end{enumerate}
\end{proposition}

\begin{proof}
For $\QH=1$: a degree-1 map $S^3\to S^2$ (Hopf map) has as its
fibres unknotted circles by the Hopf fibration.  The preimage of
any regular value is a single fibre, hence an unknot.

For $\QH=2$: the standard Hopf map $\eta:S^3\to S^2$ has Hopf
degree~1.  A degree-2 map can be constructed as $\eta\circ f$
where $f:S^3\to S^3$ is a degree-2 map.  The preimage of a regular
value under a degree-2 Hopf map consists of two circles whose
mutual linking number equals the Hopf degree:
$\QH=2\cdot\lk(\gamma_1,\gamma_2)$, so $\lk=1$ and the preimage
is the Hopf link $T(2,2)$, giving $\QH=2$.
(See Remark~\ref{P18:rem:link_convention} for the linking-number
convention.)

For $\QH=3$: the minimal-energy knot with three crossings is the
trefoil $T(2,3)$.  In the Seifert fibration of $S^3$, the trefoil
sits as the $(2,3)$-torus knot with self-linking number (framing)
$3$ in the Seifert fibered structure.  Its three crossing sites are
the structurally distinguished points identified with the three
quark vertices in Papers~XV--XVI.
\qed
\end{proof}

\begin{remark}[Linking number convention and the $\QH$ formula]
\label{P18:rem:link_convention}
Equation~\eqref{P18:eq:hopf_link} holds in the convention where the
Hopf link $T(2,2)$ has $\lk(\gamma_1,\gamma_2)=1$ and contributes
$\QH=2$ as the sum over the \emph{two-component preimage} with
total degree computed from the Seifert surface.  For a two-component
preimage $\{\gamma_1,\gamma_2\}$ with each circle winding once
around the longitude of the condensate torus, the Hopf degree is
$\QH=2\cdot({\rm winding}) = 2$, matching the two-linked-circle
picture.  The key point for what follows is not the precise formula
but the topological fact: $\QH=2$ preimage is a two-component Hopf
link, $\QH=3$ preimage is a single knotted circle (trefoil).
\end{remark}

\subsection{The satellite construction: from Hopf link to trefoil}
\label{P18:sec:satellite}

The step $\QH=2\to\QH=3$ in the condensate topology corresponds
to the satellite knot construction in classical knot theory.

\begin{definition}[Satellite knot]
\label{P18:def:satellite}
Let $K$ be a knot in $S^3$ and $P$ a knot in the solid torus
$D^2\times S^1$.  The \emph{satellite knot} $S(K,P)$ is obtained
by embedding $D^2\times S^1$ tubularly around $K$ and taking the
image of $P$ under this embedding.
\end{definition}

\begin{proposition}[Trefoil as satellite of the Hopf link]
\label{P18:prop:satellite}
The trefoil $T(2,3)$ is the satellite of the unknot $U$ with
pattern the $(2,3)$-cable, equivalently the closure of the braid
$\sigma_1^3\in B_2$ after cabling.  More precisely:
\begin{enumerate}[noitemsep,label=(\roman*)]
\item The Hopf link $T(2,2)$ contains a natural solid-torus neighbourhood
  $N(\gamma_1)$ of its first component $\gamma_1$.
\item Threading the second component $\gamma_2$ (the $\QH=1$ unknot)
  through $N(\gamma_1)$ with the correct framing converts the
  two-component link $T(2,2)\cup\gamma_2$ into the single-component
  trefoil $T(2,3)$.
\item The framing required is the $(2,3)$-cable framing: the threaded
  circle must wind $2$ times around the longitude of $N(\gamma_1)$
  and $3$ times around the meridian, which is precisely the torus knot
  winding condition $T(2,3)$.
\end{enumerate}
\end{proposition}

\begin{proof}
This is the standard cable/satellite construction for torus knots.
The $(2,3)$-cable of the unknot $U$ inside $N(U)\cong D^2\times S^1$
is the torus knot $T(2,3)$.  The connection to $T(2,2)$ follows from
the fact that $T(2,2)=L(2)$, the standard Hopf link, and its first
component bounds a Seifert disc whose boundary can be used to guide
the cabling.  See~\cite{Rolfsen1976}, Chapter~4.
\qed
\end{proof}

\begin{remark}[Physical meaning of the satellite construction]
\label{P18:rem:satellite_physical}
In the condensate, the satellite construction says: the $\QH=3$
trefoil configuration is the unique topological completion of the
$\QH=2$ Hopf-link configuration obtained by threading one additional
circle (the $\QH=1$ unknot) through the torus with the $(2,3)$-cable
framing.  This framing is not a free choice: it is the
minimal-action framing in the $\QH=3$ sector, dictated by the
same BPS saddle condition (Paper~I~\cite{Paper1}, Theorem~3.4) that
fixes $\lambda_{\rm eff}^*=\vp^6/2^{1/3}$.  The three crossing
sites of the resulting trefoil are the energetically preferred loci
where the cabling geometry meets the condensate curvature maximum.
\end{remark}

%══════════════════════════════════════════════════════════════════════════════
\section{Topological Reaction Rules}
\label{P18:sec:reaction_rules}
%══════════════════════════════════════════════════════════════════════════════

\subsection{Jones polynomial arithmetic as a selection rule}
\label{P18:sec:skein_rules}

The Jones polynomial at $q_5=e^{2\pi i/5}$ obeys the Skein relation:
\begin{equation}
  q_5^{-1}\Jones_{L_+} - q_5\Jones_{L_-}
  = (q_5^{1/2}-q_5^{-1/2})\Jones_{L_0},
  \label{P18:eq:skein}
\end{equation}
where $L_+$, $L_-$, $L_0$ differ at a single crossing.
At the root $q_5=e^{2\pi i/5}$, the coefficient
$q_5^{1/2}-q_5^{-1/2}=2i\sin(\pi/5)=i\sqrt{(5-\sqrt5)/2}$ is
irrational, so the Skein relation does not in general preserve
rationality.  Rationality is a special property of specific links
at this root.

\begin{definition}[Rationality class]
\label{P18:def:rational_class}
A satellite configuration in the twist family of
Proposition~\ref{P18:prop:satellite_twist_family} is in the
\emph{rational class} at the WZW root $q_5$ if its satellite Jones
invariant $\Jones_{\rm sat}(q_5)\in\mathbb{Q}$; otherwise it is in
the \emph{irrational class}.  By
Proposition~\ref{P18:prop:satellite_twist_family}, fused
configurations (genuine torus knots, odd $n$) are irrational and
unfused configurations (even $n$, cables not yet merged) are rational.
For the sectors of this framework: $\QH=2$ (Hopf link) has
$|\Jones_{T(2,2)}(q_5)|=1/\vp$ (irrational); $\QH=3$ (trefoil,
satellite at $n=3$) has $|\Jones_{\rm sat}(q_5)|=1/\vp$ (irrational);
and the unfused configuration at $n=4$ has
$\Jones_{\rm sat}^{(4)}(q_5)=1$ (rational).
\end{definition}

\begin{theorem}[Jones polynomial selection rule]
\label{P18:thm:jones_rule}
At the WZW root $q_5=e^{2\pi i/5}$, topological reactions
in the density-feedback condensate obey the following constraints:
\begin{enumerate}[noitemsep,label=(\roman*)]
\item \textbf{Charge conservation}: $\sum_i Q_{H,i}$ is conserved.
\item \textbf{Rationality factorisation}: If the initial state is in the
  irrational class and the final state contains a rational-class
  component (in the sense of Definition~\ref{P18:def:rational_class}),
  at least one other final-state component must be irrational.
\end{enumerate}
\end{theorem}

\begin{proof}
(i) follows from homotopy invariance of the Hopf degree.

(ii): The irrational content ($\vp$-factors) of the satellite Jones
product is carried by the fused (irrational-class) components.
A rational-class component contributes a rational factor; for
the total product to remain consistent, the remaining irrational
content must balance.  The disjoint-union factor
$-(q_5^{1/2}+q_5^{-1/2})=-\vp$ is irrational, so multi-component
states carry irrational factors that must be tracked across reactions.
Paper~XIX~\cite{Paper19} (Theorem~7.3) %\ref{P19:thm:unit_reaction_rule}
upgrades this to a unit-equivalence statement in $\mathbb{Z}[\zeta_5]$,
providing the precise bookkeeping.
\qed
\end{proof}

\begin{remark}[Proof status of Theorem~\ref{P18:thm:jones_rule}]
Theorem~\ref{P18:thm:jones_rule}(ii) is stated as a selection
rule whose precise statement requires specifying the Jones polynomial
of multi-component configurations in the condensate, which involves
a choice of disjoint-union formula that is fixed here by the Kauffman
bracket normalisation of Paper~XVII~\cite{Paper17},
Remark~11.4%\ref{P17:rem:kauffman}
.  A complete proof requires extending
the knot-invariant framework of Papers~XV--XVI to multi-component
reactions, which is the programme for Papers~XIX--XX.  We state it
as a theorem here because the structure is exact; the gap is in the
formal bookkeeping for multi-component states.
\end{remark}

%══════════════════════════════════════════════════════════════════════════════
\section{Forced Pion Emission}
\label{P18:sec:pion}
%══════════════════════════════════════════════════════════════════════════════

\begin{theorem}[Pion emission is topologically forced]
\label{P18:thm:pion_forced}
The reaction
\begin{equation}
  \QH=2 \;+\; \QH=2 \;\longrightarrow\; \QH=3 \;+\; \QH=1
  \label{P18:eq:reaction_pion}
\end{equation}
is the unique minimal topological reaction that:
\begin{enumerate}[noitemsep,label=(\roman*)]
\item conserves Hopf charge ($2+2=3+1$);
\item converts two irrational-class objects into one rational-class
  object (the $\QH=3$ trefoil) with one irrational-class companion;
\item minimises the number of components in the final state.
\end{enumerate}
The $\QH=1$ product is identified with the pion ($\pi^\pm$ or
$\pi^0$) of the physical baryon sector.  Pion emission in baryon
formation is therefore a topological necessity.
\end{theorem}

\begin{proof}
Charge conservation requires the final state to have total Hopf
charge 4.  The minimal decomposition into sectors with $\QH\geq1$
that includes a $\QH=3$ object is $4=3+1$.  Alternatives are
$4=3+1$, $4=2+2$ (no $\QH=3$), $4=2+1+1$, $4=1+1+1+1$ (all $\QH=1$).
Only $3+1$ satisfies both the charge conservation and the requirement
of producing a confined ($\QH=3$) baryon.  The $\QH=1$ remainder is
therefore uniquely forced.

At the preimage level: two $\QH=2$ Hopf links (four circles total)
must reconfigure into the trefoil preimage.  The trefoil is a
single-component knot.  Of the four initial circles, three must
fuse into the trefoil via the satellite construction
(Proposition~\ref{P18:prop:satellite}); the fourth remaining circle
becomes the $\QH=1$ unknot spectator.  This fourth circle cannot
be absent: a four-circle configuration cannot collapse to a single
knot without producing a residual, as the linking number is conserved
modulo the reaction topology.
\qed
\end{proof}

\begin{corollary}[Pion as topological spectator]
\label{P18:cor:pion_spectator}
In the reaction~\eqref{P18:eq:reaction_pion}, the $\QH=1$ pion is the
preimage circle that does not participate in forming the trefoil.
Its internal dynamics are governed by $\SU(2)_1$ WZW theory with
$V^*=1$ and Majorana statistics (Paper~XVII~\cite{Paper17},
Corollary~9.13%\ref{P17:cor:majorana}
).  The pion is the lightest
topological object and the necessary companion of every
topological baryon-formation event.
\end{corollary}

\begin{remark}[$\QH=2+\QH=1$ reaction]
\label{P18:rem:qh2_qh1}
The complementary reaction
$\QH=2+\QH=1\to\QH=3$ is the satellite construction directly:
the $\QH=1$ unknot threads through the $\QH=2$ Hopf link with
the $(2,3)$-cable framing to produce the trefoil.  This reaction
has no spectator, since all three preimage circles are consumed
in forming the trefoil.  Energetically this is the simpler
baryon-formation channel; it corresponds physically to the
absorption of a pion by a lepton-sector object to produce a baryon,
i.e.\ the topological analogue of $e^-+\pi^+\to p$.
\end{remark}

\begin{remark}[Higher-charge incoming objects and minimality of the listed channels]
\label{P18:rem:higher_charge_incoming}
Theorem~\ref{P18:thm:pion_forced} and Remark~\ref{P18:rem:qh2_qh1} list the
\emph{minimal} charge-conserving production channels for $\QH=3$.
A natural question is whether the incoming perturbation could carry
$\QH>1$: topological charge conservation requires only that the
\emph{total} incoming charge sum to~$3$, so the full set of allowed
reactions includes, for example,
\[
  \QH=3 + \QH=0 \;\to\; \QH=3
  \quad\text{(trivial; } \QH=0 \text{ is the condensate ground state)},
\]
as well as in principle any decomposition $n_1+n_2+\cdots=3$ with
$n_i\geq1$.  In particular a $\QH=4$ or higher object could in
principle collide with a $\QH=2$ target and deposit charge.

The two channels listed are the physically minimal cases for three
independent reasons.  First, they involve only sectors that the
framework has established as stable or well-defined: $\QH=1$
(neutrino/pion, stable, Paper~XVII~\cite{Paper17}) and $\QH=2$
(charged lepton, stable, Paper~I~\cite{Paper1}).  A $\QH\geq4$
incoming object would itself be transient --- the framework
establishes no stability or metastability for $\QH\geq4$ --- and
whether a transient-on-transient interaction can produce a
stable $\QH=3$ before both decay is a dynamical question outside
the current scope.  Second, the listed channels are the
lowest-energy threshold reactions: higher incoming $\QH$ requires
higher total production energy, so they are the first channels to
open as collision energy increases.  Third, the
$\QH=2+\QH=2\to\QH=3+\QH=1$ channel is already the one identified
with real hadronic processes (the topological analogue of
$e^-+e^-\to p+\bar\nu$), so the minimal cases are also the
physically realised ones.

A related point is that charge conservation plus energy minimisation
also determines the final state: a $\QH=2+\QH=2$ collision could
in principle produce $\QH=4$ (conserving charge) rather than
$\QH=3+\QH=1$.  The $\QH=4$ final state is disfavoured for the same
reason a $\QH\geq4$ incoming object is: no stable $\QH=4$
configuration is established by the framework, and the energetic cost
of a higher-charge object exceeds that of the $3+1$ split.  The
selection principle is therefore: topological charge conservation
permits any channel whose charges sum correctly; energy minimisation
and stability of the final-state objects together select the minimal
split into established-stable sectors.
\end{remark}

\begin{remark}[Kinematic path of the $\QH=2\to\QH=3$ transition]
\label{P18:rem:crossing_path}
The satellite construction of Proposition~\ref{P18:prop:satellite} names
the topological endpoint of the $\QH=2+\QH=1\to\QH=3$ reaction but
does not describe the \emph{geometric path} from start to finish.
The following three-step sequence provides that path.

\textbf{Preliminary: the $\QH=2$ Hopfion is two tubes, not one.}
The density-feedback stabilised $\QH=2$ Hopfion is built on the
Hopf link $T(2,2)$: two unknotted circles, each an independent tube
of condensate field, linked once around each other.  This is not a
single tube reconfiguring; it is two distinct tube components
$\gamma_1$ (outer) and $\gamma_2$ (inner) from the outset.  Because
the density-feedback functional $K_{\rm fb}+\vp^6 J_4$ is
rotationally symmetric around the torus axis, the two tubes are
energetically equivalent: there is no preferred point along either
component where a perturbation is more likely to act.  The
``splitting point'' that the construction assigns at $t=\pi$ is a
calculational convenience, not a physical feature; any toroidal
angle would give the same field by symmetry.

\textbf{Step~1 (writhe injection, inner tube deforms).}  The
incoming $\QH=1$ perturbation acts on the inner component $\gamma_2$,
injecting one writhe unit and deforming its path from a planar
unknotted circle into a figure-eight-like intermediate.  This is a
\emph{geometric}, not topological, change: the linking number of
$\gamma_2$ with $\gamma_1$ is unchanged, and no crossing has been
committed.  The outer component $\gamma_1$ and its neighbourhood
$N(\gamma_1)$ remain approximately unperturbed at this stage.

\textbf{Step~2 (threading begins, inner approaches outer core).}
The deformed inner component $\gamma_2$ begins to approach the
interior of the outer component's solid-torus neighbourhood
$N(\gamma_1)$.  The approach is \emph{forced to be asymmetric}
by the $(2,3)$-cable framing constraint of
Proposition~\ref{P18:prop:satellite}: $\gamma_2$ must wind twice around
the longitude and three times around the meridian of $N(\gamma_1)$.
A symmetric approach --- $\gamma_2$ approaching the core of
$N(\gamma_1)$ from directly opposite --- would instead produce a
doubly-wound configuration, which is a different topological sector
($\QH=2$ doubly wound, energetically penalised by the
density-feedback suppression).  The asymmetric threading is
therefore not a contingent feature of any particular perturbation:
it is the only geometric path consistent with the framing
constraint.  As threading proceeds, $\gamma_2$'s own loop tightens
to accommodate the winding, with the tube topology of both
components preserved throughout.

\textbf{Step~3 (single-site threading generates all three crossings).}
The three crossings of the trefoil $T(2,3)$ do not appear as three
discrete sequential events with the threading choosing among
alternatives.  Once $\gamma_2$ has threaded through $N(\gamma_1)$
with the $(2,3)$-cable framing of Proposition~\ref{P18:prop:satellite},
all three crossings of $T(2,3)$ are geometrically determined
simultaneously as continuous consequences of the completed cable
winding.  In the braid group language, the trefoil is the closure
of $\sigma_1^3\in B_2$: the exponent $3$ counts the crossings of the
resulting knot, but the threading event that produces them is
\emph{one} continuous geometric process, not three discrete choices.
The ``unique valid crossing'' of the transition is therefore the
\emph{single perturbation site} where the threading of $\gamma_2$
into $N(\gamma_1)$ begins --- and this site is unique in the sense
that its location is set by the initial $\QH=1$ perturbation, but
once the threading commences the three crossings of the completed
$T(2,3)$ are fixed by the cable framing, not chosen individually.

This picture is directly visible in the $T(2,q)$ construction sweep
of Figure~\ref{P18:fig:qpol_sweep}: as $q_{\rm pol}$ moves continuously
from $2.0$ toward $3.0$ (with $\alpha_{\rm wind}=0$ throughout), the
single perturbation site of $q=2.3$ evolves continuously into the
three-lobe, three-crossing structure of $q\to3$, without ever
requiring a discrete ``choice'' between multiple candidate crossing
sites.  The rubber-band picture of a single tube reconfiguring, in
which one might imagine multiple candidate crossings competing to
be committed, is misleading here: the $\QH=2$ Hopfion is two tubes
throughout, and the transition is one threading event, not three
crossing events.

\textbf{The crossing is the decay site.}  Step~3 completes the
$\QH=3$ topology, but the crossing vertex that does so is
simultaneously the site of maximum $S_{\rm eff}$ and minimum
$\rho_{J_4}$ (Proposition~\ref{P18:prop:anticorrelation}).  The
anti-correlation is causal: the density-feedback suppression is
strongest at the crossings (maximum curvature $\kappa=0.399$),
which depletes the topological flux there to its minimum.  The
crossing is the decay site not because suppression is weak but
because the field is thinnest: the topological bottleneck where
suppression is strongest is also where the structure is most
vulnerable to reversal, since the minimum $\rho_{J_4}$ means
the least energy is locally available to sustain the crossing.
The consequence is that decay begins at formation.  A candidate
explanation --- not established, offered here as a conjecture
consistent with the abstract's own hedging of this same point --- is
that the crossing vertex has no stable irrational algebraic anchor
because the formation event passes through the unfused intermediate
configuration ($n=4$ in the twist family of
Proposition~\ref{P18:prop:satellite_twist_family}), whose satellite
Jones invariant is rational ($\Jones_{\mathrm{sat}}^{(4)}(q_5)=1$,
equation~\eqref{P18:eq:sat_n4}), before fusing to the irrational
$n=3$ value ($|\Jones_{\mathrm{sat}}^{(3)}(q_5)|=1/\vp$,
Theorem~\ref{P18:thm:trefoil_satellite_jones}).  On this conjectural
reading, a rational invariant provides no irrational BPS fixed point
in the profile space, offering one possible account of why the
$\QH=3$ energy landscape has a flat plateau
(Remark~\ref{P18:rem:saddle_difficulty}) rather than a unique sharp
minimum; whether the formation path genuinely passes through $n=4$ is
not established (Section~\ref{P18:sec:jones}).  What \emph{is}
established, independent of this conjecture, is that the flat plateau
itself is an observed numerical fact (Papers~XV/XVI gradient-flow
searches) and that the satellite invariant is irrational at $n=3$: on
either reading, the formation event reaches $n=3$ with the crossing
already trying to reverse.
What sustains the $\QH=3$ configuration transiently is not an
energy barrier but the incoming kinetic energy from the $\QH=2+\QH=1$
reaction, which must be redistributed through the cascade of
Remark~\ref{P18:rem:implode_explode} before the crossing can reverse.
The baryon is metastable in the dynamical sense (it persists while
the kinetic energy is being shed) not the energetic sense (there is
no potential-energy barrier preventing decay).  This distinction
is the correct reading of the ``flat-bottomed $\QH=3$ energy
landscape'' of Remark~\ref{P18:rem:saddle_difficulty}: flat means no
barrier, not no configuration.
\end{remark}

\begin{remark}[Two sub-thresholds and the virtual baryon regime]
\label{P18:rem:two_thresholds}
Remark~\ref{P18:rem:crossing_path} describes the geometry of a successful
$\QH=2\to\QH=3$ transition.  The energy threshold of
Remark~\ref{P18:rem:production_conditions} controls whether that
transition occurs at all --- but the threshold is not a single
number.  It has internal structure: two distinct sub-thresholds
separate three physically distinct regimes.

\textbf{Sub-threshold I (commitment threshold).}  This is the minimum
kinetic energy required to drive the fold of Step~2 and the crossing
of Step~3 in Remark~\ref{P18:rem:crossing_path} against the restoring
force of the density-feedback suppression.  Below this threshold the
reaction fails at Step~2 or Step~3: the fold reverts before the
cable framing is achieved, or the crossing forms with the wrong
handedness and immediately unwinds.  In either case the field returns
elastically to the $\QH=2$ sector.  This is the dynamical mechanism
of confinement in the framework: below the commitment threshold the
condensate actively restores $\QH=2$, not merely by forbidding free
quarks by fiat but by reversing the geometric sequence that would
have produced them.

\textbf{Sub-threshold II (ejection threshold).}  This is the minimum
kinetic energy required for the implosion--explosion cascade of
Remark~\ref{P18:rem:implode_explode} to run to completion before the
committed crossing reverses.  The crossing is immediately trying to
reverse (Remark~\ref{P18:rem:crossing_path}: decay begins at formation),
so the cascade --- lobe implosion, inward energy redistribution,
redirection through the crossing bottleneck (where $S_{\rm eff}$
is maximal and $\rho_{J_4}$ is minimal ---
Proposition~\ref{P18:prop:anticorrelation}),
geometric separation of the ejected $\QH=1$ object --- must complete
before the crossing unwinds.  Once the daughter $\QH=1$ object has
separated far enough that reconnection is geometrically impossible,
ejection is irreversible.  The ejection threshold is therefore
strictly higher than the commitment threshold, because it requires
not only committing the crossing but sustaining it through the full
cascade timescale.

\textbf{The virtual baryon regime.}  Between the two sub-thresholds
lies a regime in which the $\QH=3$ crossing is committed but the
cascade cannot complete before reversal.  The field momentarily
achieves $\QH=3$ topology, then the crossing unwinds and the system
returns to $\QH=2$ without ejecting a $\QH=1$ daughter.  This is a
\emph{virtual} $\QH=3$ fluctuation: topologically real while it
exists, but producing no observable final-state quark.  Whether this
regime is physically wide (the two sub-thresholds are well-separated)
or narrow (the cascade is fast relative to the reversal timescale and
the sub-thresholds nearly coincide) is an open question, addressable
by the kick-strength sweep of the dynamical integrator described in
the companion numerical programme (sweeping \texttt{kick\_strength}
from the $\QH=2$ starting state and watching $J_4$ for a transient
excursion into $\QH=3$ territory followed by return rather than
sustained ejection).

\textbf{Summary.}  Any perturbation above the ejection threshold
commits the $\QH=3$ crossing and sustains it through the full
cascade, ejecting a $\QH=1$ daughter irreversibly.  Any perturbation
between the two sub-thresholds commits the crossing but fails to
sustain it, producing a virtual $\QH=3$ fluctuation with no
observable quark emission.  Any perturbation below the commitment
threshold restores $\QH=2$ elastically.  The energy threshold of
Remark~\ref{P18:rem:production_conditions} should be read as the
\emph{ejection} threshold (the physically observable one); the
commitment threshold and the width of the virtual-baryon regime are
the quantitative open questions that the dynamic integrator sweep
would resolve.
\end{remark}

\begin{figure}[htp]
    \centering
    \begin{subfigure}[T]{\linewidth}
        \begin{minipage}[T]{0.3\textwidth}
            \includegraphics[width=0.95\linewidth]{n_initial_q2.00_a0.00.png}
        \end{minipage}\begin{minipage}[T]{0.7\textwidth}
            \caption{$q=2.0$, $\alpha=0$: the $T(2,2)$ Hopf-link
            construction with $\QH=2$ field winding. The two-tube structure
            forms a clean ring with smooth color gradient of the two linked
            circles $\gamma_1$ (outer) and $\gamma_2$ (inner), each an
            independent tube of condensate field. The slight visible
            strain is an artefact of the two-strand construction
            ansatz (the density-feedback stabilised $\QH=2$ lepton
            vacuum is the gradient-flow saddle at $C_2^*=3.4318$, not
            this initial construction).}
        \end{minipage}
        \label{P18:fig:qpol_q20}
    \end{subfigure}
    \begin{subfigure}[T]{\linewidth}
        \begin{minipage}[T]{0.3\textwidth}
            \includegraphics[width=0.95\linewidth]{n_initial_q2.30_a0.00.png}
        \end{minipage}\begin{minipage}[T]{0.7\textwidth}
            \caption{$q=2.3$, $\alpha=0$: the threading event has begun
            at a single site (visible on the left, approximately
            $9$~o'clock), where the two arc-halves of the
            construction can no longer be stitched together smoothly
            because the $T(2,q)$ curve no longer closes in $[0,\pi)$
            for $q>2$.  The inner tube $\gamma_2$ has begun to
            deform toward and thread through $N(\gamma_1)$ at this
            single location; the topology is still $\QH=2$
            ($\alpha=0$, no winding commitment).  This site is the
            geometric precursor of the trefoil threading of
            Remark~\ref{P18:rem:crossing_path}.}
        \end{minipage}
        \label{P18:fig:qpol_q23}
    \end{subfigure}
    \begin{subfigure}[T]{\linewidth}
        \begin{minipage}[T]{0.3\textwidth}
            \includegraphics[width=0.95\linewidth]{n_initial_q2.50_a0.00.png}
        \end{minipage}\begin{minipage}[T]{0.7\textwidth}
            \caption{$q=2.5$, $\alpha=0$: the threading of $\gamma_2$
            through $N(\gamma_1)$ is well underway.  The inner tube has
            deformed substantially into the outer tube's neighbourhood,
            and the three-lobe structure of the incipient trefoil is
            already visible even though the topology is still $\QH=2$
            ($\alpha=0$).  The construction curve does not close in
            $[0,2\pi)$ at non-integer $q$; consequently the field is
            somewhat disorganised across the whole ring, but the
            dominant geometric feature remains the single threading
            event initiated at $q\approx2.3$.}
        \end{minipage}
        \label{P18:fig:qpol_q25}
    \end{subfigure}
    \begin{subfigure}[T]{\linewidth}
        \begin{minipage}[T]{0.3\textwidth}
            \includegraphics[width=0.95\linewidth]{n_initial_q2.70_a0.00.png}
        \end{minipage}\begin{minipage}[T]{0.7\textwidth}
            \caption{$q=2.7$, $\alpha=0$: At the field-orientation
            discontinuity at approximately 9 o'clock, the inner component
            $\gamma_2$ is approaching the interior of $N(\gamma_1)$: the
            three-lobe structure of the trefoil is now well defined and
            geometrically present while still under the $\mathcal{Q}_H=2$
            topology. The strain site of $q=2.5$ is beginning to collapse 
            toward a tube closure site as $q\to 3$.}
        \end{minipage}
        \label{P18:fig:qpol_q27}
    \end{subfigure}
    \begin{subfigure}[T]{\linewidth}
        \begin{minipage}[T]{0.3\textwidth}
            \includegraphics[width=0.95\linewidth]{n_initial_q2.90_a0.00.png}
        \end{minipage}\begin{minipage}[T]{0.7\textwidth}
            \caption{$q=2.9$, $\alpha=0$: near-trefoil geometry.  The
            construction curve is close to $T(2,3)$ but not yet closed
            in $[0,2\pi)$ (closure gap $\approx0.56$ in coordinate
            units).  The two-tube structure of $q=2.0$ has consolidated
            into approximately a single continuous tube with three lobes
            --- the characteristic trefoil geometry --- via the single
            threading event initiated at $q\approx2.3$.  At $q=3$ the
            curve closes and the configuration completes to a proper
            $T(2,3)$; the topological commitment
            ($\QH=2\to\QH=3$) is controlled separately by
            $\alpha_{\rm wind}$.}
        \end{minipage}
        \label{P18:fig:qpol_q29}
    \end{subfigure}
    \caption{Geometric path of the $\QH=2\to\QH=3$ transition,
    visualised via the $T(2,q)$ construction family at fixed
    $\QH=2$ field winding ($\alpha_{\rm wind}=0$, so $\Phi=\chi+2\tau$
    throughout), viewed in directional-arrow mode.
    Arrow direction encodes $\mathbf{n}(\mathbf{x})$;
    colour encodes $\mathbf{n}$ direction on $S^2$.
    All panels: $N=64$, $h=0.175$, $R_0=3$, $r_0=0.874$, $C^*=2.5062$,
    produced by \texttt{hopf\_link\_construction_v2.py} with
    \texttt{--verify\_only} (no gradient flow).
    The parameter $q_{\rm pol}\in[2,3]$ controls the curve geometry
    only; the field topology is $\QH=2$ at all five panels.
    The sequence shows the $\QH=2$ Hopfion as two linked tubes at
    $q=2.0$ (Panel~\subref{fig:qpol_q20}), the single-site threading
    event initiating at $q=2.3$ (Panel~\subref{fig:qpol_q23},
    approximately 9~o'clock), progressing at $q=2.5$
    (Panel~\subref{fig:qpol_q25}), and consolidating into the
    near-trefoil geometry at $q=2.9$ (Panel~\subref{fig:qpol_q29}).
    The transition proceeds via a single continuous threading event,
    not via multiple discrete crossings competing to be chosen; the
    three crossings of the completed $T(2,3)$ emerge simultaneously
    as continuous consequences of the completed cable framing
    (Remark~\ref{P18:rem:crossing_path} Step~3).  The topological
    commitment ($\alpha_{\rm wind}=0\to1$) is a separate parameter
    that is zero throughout this figure.  Note that $q=2.0$
    carries a slight construction strain (the true density-feedback
    stabilised $\QH=2$ vacuum is the gradient-flow saddle at
    $C_2^*=3.4318$; the construction shown here is the initial
    two-strand ansatz, not the converged saddle field).}
    \label{P18:fig:qpol_sweep}
\end{figure}

%──────────────────────────────────────────────────────────────────────────────
\subsection{The $T(2,q)$ geometric interpolation: geometry without topology}
\label{P18:sec:twoq_interpolation}
%──────────────────────────────────────────────────────────────────────────────

Figure~\ref{P18:fig:qpol_sweep} makes visible something that the algebraic
description of Remarks~\ref{P18:rem:crossing_path} and~\ref{P18:rem:two_thresholds}
can only state abstractly: the geometric threading site that
eventually commits $\QH=3$ topology is already present in the field
before any topological commitment has been made, and it is a
\emph{single} site, not a set of alternatives.  The five panels
sweep the construction curve from $T(2,2)$ to $T(2,2.9)$ while
holding the field winding fixed at $\alpha_{\rm wind}=0$ (pure
$\QH=2$, $\Phi=\chi+2\tau$ throughout), so the field topology is
$\QH=2$ at every panel.  The only thing changing is the curve
geometry, controlled by $q_{\rm pol}$.

\begin{observation}[Geometric precursor at $q=2.3$]
\label{P18:obs:precursor}
At $q_{\rm pol}=2.3$, the directional-arrow field already displays a
sharp field-orientation discontinuity at one site on the ring (left,
approximately 9~o'clock in (Figure~\ref{P18:fig:qpol_sweep}\subref{fig:qpol_q23})).
This is not a topological event: the Whitehead $\QH$ of the field remains
$2$ at this value, and the linking number of $\gamma_2$ with $\gamma_1$ is
unchanged.  It is a geometric event: the $T(2,q)$ curve no longer
closes in $[0,\pi)$ for $q>2$, so the two arc-halves can no longer
be joined smoothly with the same Bishop frame, and the construction
field records the mismatch as a strain site.

The location of this strain site is the geometric precursor of the
single threading event of Step~2 (Remark~\ref{P18:rem:crossing_path}).
The field is showing, at $q=2.3$ and topology $\QH=2$, exactly where
the eventual topological commitment will occur if $\alpha_{\rm wind}$
is increased toward~$1$: this single site, not a set of independent
candidate crossings, is where the threading of $\gamma_2$ into
$N(\gamma_1)$ begins.  Geometry precedes topology: the threading
site is identified before the topological commitment is made.
\end{observation}

The progression across the five panels makes the geometric structure
of Remark~\ref{P18:rem:crossing_path} visible as a continuous
deformation.

\textbf{$q=2.0$ (Figure~\ref{P18:fig:qpol_sweep}\subref{fig:qpol_q20}).}  The $T(2,2)$
Hopf-link construction with $\QH=2$ field winding.  The two-tube
structure is directly visible: two distinct linked tubes
$\gamma_1$ (outer) and $\gamma_2$ (inner).  The slight visible
strain is an artefact of the two-strand ansatz (the true
density-feedback stabilised $\QH=2$ vacuum is the gradient-flow
saddle at $C_2^*=3.4318$, not this initial field).  The two tubes
are energetically equivalent at this stage: there is no preferred
point along either component where a perturbation would
preferentially act, consistent with the rotational symmetry of
$K_{\rm fb}+\vp^6 J_4$.

\textbf{$q=2.3$ (Figure~\ref{P18:fig:qpol_sweep}\subref{fig:qpol_q23}).}  The inner component
$\gamma_2$ has deformed under the incoming $\QH=1$ perturbation: one
writhe unit has been injected, and the threading of $\gamma_2$ into
$N(\gamma_1)$ has begun at a single site (visible at approximately
9~o'clock in the panel).  This corresponds to Step~1 of
Remark~\ref{P18:rem:crossing_path} --- writhe injection, inner tube
deforms --- with the strain site identifying the figure-eight
intermediate expressed in the field.  The outer component $\gamma_1$
and $N(\gamma_1)$ remain approximately unperturbed away from the
threading site.

\textbf{$q=2.5$ (Figure~\ref{P18:fig:qpol_sweep}\subref{fig:qpol_q25}).}  The threading of
$\gamma_2$ through $N(\gamma_1)$ is well underway.  The inner tube
has deformed substantially into the outer tube's neighbourhood, and
the three-lobe structure of the incipient trefoil is already visible
even though the topology is still $\QH=2$.  Because the $T(2,q)$
curve does not close in $[0,2\pi)$ for non-integer $q$, the field
across the ring is somewhat disorganised, but the dominant geometric
feature remains the single threading event initiated at $q\approx2.3$
--- not a set of independent competing crossings.

\textbf{$q=2.7$--$2.9$ (Figure~\ref{P18:fig:qpol_sweep}[\subref{fig:qpol_q27},
\subref{fig:qpol_q29}]).}  The threading
has consolidated: what began at $q=2.3$ as a strain concentrated at a
single site has, by $q=2.9$, become the near-complete trefoil
geometry.  The two-tube structure of $q=2.0$ has consolidated into
approximately a single continuous tube with three lobes, with the
$T(2,q)$ curve close to $T(2,3)$ (closure gap $\approx0.56$ in
coordinate units) and $\gamma_2$'s loop tightened substantially to
accommodate the cable winding.  All three trefoil crossings are now
geometrically present as continuous consequences of the completed
single threading event, not as three independently chosen events
--- a physical distinction that the T(2,q) sweep makes visible in a
way the algebraic braid-group description
($\sigma_1^3\in B_2$) cannot.

\textbf{What $\alpha_{\rm wind}=0$ shows throughout.}  The separation
of $q_{\rm pol}$ and $\alpha_{\rm wind}$ into independent parameters
exposes what the full $T(2,3)$ construction conflates.  The entire
$q=2.0\to2.9$ progression is purely geometric: the topology is $\QH=2$
at every step, and no topological commitment has occurred.
Once $\alpha_{\rm wind}$ is increased toward~$1$ at fixed
$q_{\rm pol}\approx2.9$, the field winding transitions from $\QH=2$ to
$\QH=3$ and the condensate enters the baryon sector.  The single
threading site of the $\alpha_{\rm wind}=0$ sweep and the three
crossings of the fully committed $T(2,3)$ are the same geometric
event, viewed at different values of the topological commitment
parameter: the sweep at $\alpha_{\rm wind}=0$ shows the threading
without commitment; the completed $T(2,3)$ at
$\alpha_{\rm wind}=1$ shows the same event with the topology
committed.


\begin{remark}[Production conditions: a two-threshold structure]
\label{P18:rem:production_conditions}
Theorem~\ref{P18:thm:pion_forced} and Remark~\ref{P18:rem:qh2_qh1} establish
\emph{which} reactions are topologically permitted; they say nothing
about \emph{when} a $\QH=3$ object actually forms.  Two independent
physical conditions must both be satisfied, neither of which is
addressed by the topological selection rule alone.

\textbf{(i) Energy threshold.}  Writing $E(\QH{=}n)$ for the physical
rest energy of a sector-$n$ object on a common scale, the reactions
require
\[\Delta E_{2+2} = E(\QH{=}3)+E(\QH{=}1)-2E(\QH{=}2) \;>\;0\]
\[\qquad\text{or}\qquad\]
\[\Delta E_{2+1} = E(\QH{=}3)-E(\QH{=}2)-E(\QH{=}1) \;>\;0\]
to be supplied as kinetic energy of the incoming objects.  With the
physical identification $\QH=3\leftrightarrow$ baryon,
$\QH=2\leftrightarrow$ charged lepton, $\QH=1\leftrightarrow$
neutrino/pion (Paper~XVII~\cite{Paper17}), $E(\QH{=}3)$ is set by the
baryon mass scale, which exceeds the charged-lepton and pion rest
masses by two to three orders of magnitude on the physical scale.
Both reactions are therefore kinematically closed at rest: they
require the colliding objects to carry substantial kinetic energy,
not merely to exist.  This is the topological-sector statement of an
ordinary threshold reaction, structurally identical to
$e^-+\pi^+\to p$ (the real-world analogue already noted in
Remark~\ref{P18:rem:qh2_qh1}) requiring the pion or electron to be
relativistic.  The dimensionless condensate-normalisation ratio
$E_{\rm fb}^{(3)}/E_{\rm fb}^{(2)}=N_3^2/N_2^2=9/100$ (Paper~XV,
\S\,``The lepton--quark hierarchy'') is a statement about the two
sectors'
\emph{own} saddle energy scales and should not be read as a
reaction-threshold estimate: $\Delta E$ requires the physical mass
formulas of Paper~XVII, not this dimensionless ratio.

\textbf{(ii) Rate/density threshold.}  Even when $\Delta E>0$ is
locally available as kinetic energy, the reaction occurs at a rate
proportional to the local number density of the colliding species:
$\Gamma_{2+2}\propto n_{\QH=2}^2\langle\sigma v\rangle$ for the
$\QH{=}2+\QH{=}2$ channel, $\Gamma_{2+1}\propto
n_{\QH=2}\,n_{\QH=1}\langle\sigma v\rangle$ for the satellite channel.
In a dilute vacuum, even a population of individually energetic
$\QH=2$ objects encounters a collision partner too rarely for
$\QH=3$ formation to proceed at any appreciable rate.

\textbf{Where both conditions hold simultaneously.}  Ordinary
laboratory and astrophysical vacuum satisfies neither condition.
Three environments satisfy both:
\begin{enumerate}[noitemsep,label=(\alph*)]
\item \emph{Collider beams}: the energy threshold is supplied
  externally (beam energy), and the density threshold is irrelevant
  for a single engineered beam-beam collision.
\item \emph{The primordial quark-gluon plasma}: thermal energy
  $kT\sim T_c$ near the QCD deconfinement temperature supplies the
  energy threshold, while the high ambient number density of the
  early universe supplies the rate threshold; both conditions hold
  together automatically as the plasma cools through $T_c$, which is
  the standard picture of baryon formation in the early universe.
\item \emph{Compact-object extremes} (neutron star cores, merger
  events, black hole horizons): extreme density supplies the rate
  threshold, and the gravitational/thermal energy released supplies
  the kinetic threshold.
\end{enumerate}
This two-threshold structure --- not a single internal-field
quantity such as $S_{\rm eff}(\theta,\rho)$ of
Proposition~\ref{P18:prop:anticorrelation}, which describes the
\emph{internal} structural weak points of an already-formed
$\QH=3$ object rather than the conditions for its formation --- is
the correct framework-level statement of why $\QH=3$ objects are
rare in the vacuum and require special formation conditions.  A
quantitative version, computing $\Delta E_{2+2}$ and
$\Delta E_{2+1}$ directly from the mass formulas of
Paper~XVII~\cite{Paper17} rather than from the qualitative
identification above, is Open Problem~O6 below.
\end{remark}

%══════════════════════════════════════════════════════════════════════════════
\section{The $\phistar$ Prediction from $T_{2,3}$ Geometry}
\label{P18:sec:phi_star}
%══════════════════════════════════════════════════════════════════════════════

\subsection{The acoplanarity angle $\phistar$}
\label{P18:sec:phi_star_def}

Let $\rho^+\rho^-$ be a pair of $\rho$ mesons (each $\QH=2$)
produced in the rest frame of the pair from the string-breaking of
a $\QH=3$ baryon.  The \emph{acoplanarity angle} $\phistar$ is
defined as the angle between the two decay planes in the $\rho^+\rho^-$
rest frame:
\begin{equation}
  \cos\phistar = \frac{(\hat{n}_+\times\hat{p})\cdot
                       (\hat{n}_-\times\hat{p})}
                      {|\hat{n}_+\times\hat{p}|\,|\hat{n}_-\times\hat{p}|},
  \label{P18:eq:phi_star_def}
\end{equation}
where $\hat{n}_\pm$ are the unit normals to the $\rho^\pm$ decay
planes and $\hat{p}$ is the beam direction (or the $\rho^+$
momentum direction in the pair CM frame).  The distribution
$W(\phistar)$ is the main observable.

\subsection{Geometric derivation of $W(\phistar)$}
\label{P18:sec:phi_star_geom}

The trefoil $T_{2,3}$ with standard parameters $R_0=3$, $r_0$
(Paper~XVI~\cite{Paper16}) is embedded as:
\begin{equation}
  \Gamma(t) = \bigl((R_0+r_0\cos 3t)\cos 2t,\;
               (R_0+r_0\cos 3t)\sin 2t,\;
               r_0\sin 3t\bigr), \quad t\in[0,2\pi).
  \label{P18:eq:trefoil_param}
\end{equation}
The three crossing vertices occur at
$t_n = \pi/6 + 2\pi n/3$, $n=0,1,2$ (Paper~XVI,
Proposition~13.3).%\ref{P16:prop:geom_data}

\begin{theorem}[Crossing geometry of $T_{2,3}$]
\label{P18:thm:crossing_geometry}
For the parametrisation~\eqref{P18:eq:trefoil_param}:
\begin{enumerate}[noitemsep,label=(\roman*)]
\item The tangent vector $\dot\Gamma(t_n)$ lies exactly in the
  $z=0$ plane at each crossing: $\dot\Gamma_z(t_n)=3r_0\cos(3t_n)
  =3r_0\cos(\pi/2)=0$.
\item The azimuthal angles of $\dot\Gamma(t_n)$ are
  $\alpha_n = 2t_n+\pi/2 = \pi/3+2\pi n/3+\pi/2$ modulo $2\pi$,
  separated by exactly $2\pi/3$ for $n=0,1,2$.
\item Each pair of adjacent crossings is separated by exactly
  $2\pi/3$ in the azimuthal direction.
\end{enumerate}
\end{theorem}

\begin{proof}
(i): $\dot\Gamma_z(t)=3r_0\cos(3t)$.  At $t_n=\pi/6+2\pi n/3$:
$3t_n=\pi/2+2\pi n$, so $\cos(3t_n)=0$ exactly.

(ii): The azimuthal direction of the tangent is determined by
$\dot\Gamma_x$ and $\dot\Gamma_y$ at $t=t_n$, where by (i) the
$z$-component vanishes.  Computing:
$\dot\Gamma_{xy}(t_n)\propto(-\sin(2t_n),\cos(2t_n))$
(the azimuthal unit vector at angle $2t_n+\pi/2$).
For $n=0,1,2$: the angles $2t_n=\pi/3+4\pi n/3$ are
$\pi/3$, $\pi/3+4\pi/3$, $\pi/3+8\pi/3$ modulo $2\pi$,
each separated by $4\pi/3$, which is $2\pi/3$ modulo $2\pi/3$.

(iii) follows from (ii) by the $\Zthree$ symmetry of the parametrisation.
\qed
\end{proof}

\begin{theorem}[$\phistar$ distribution for baryon-initiated production]
\label{P18:thm:phi_star}
When a $\QH=3$ baryon fragments to produce a $\rho^+\rho^-$ pair
at two of its three crossing vertices, and the decay-plane
orientation of each $\rho$ meson is set by the tangent direction at
its production crossing, the $\phistar$ distribution is:
\begin{equation}
  W(\phistar) = \frac{1}{2\pi}\bigl[1 + a_3\cos(3\phistar)\bigr],
  \qquad a_1=a_2=0,
  \label{P18:eq:phi_star_dist}
\end{equation}
with $a_3\neq0$ and the distribution peaked at $\phistar=2\pi/3$
(and its images under $\Zthree$).
\end{theorem}

\begin{proof}
By Theorem~\ref{P18:thm:crossing_geometry}(iii), adjacent crossing
tangents are separated by $2\pi/3$.  The decay-plane normal of a
$\rho$ meson produced at crossing $n$ is determined by the tangent
direction $\hat\alpha_n$ at that crossing (since energy ejected
along the tangent in the $z=0$ plane defines the decay plane).  For
$\rho^+$ produced at crossing $n$ and $\rho^-$ at crossing $n+1$,
the relative angle between their decay-plane normals is
$\phistar=\alpha_{n+1}-\alpha_n=2\pi/3$ (by
Theorem~\ref{P18:thm:crossing_geometry}).

The distribution $W(\phistar)$ is a function on $[0,2\pi)$ that
must be invariant under the $\Zthree$ symmetry ($\phistar\to
\phistar+2\pi/3$).  The Fourier decomposition
$W(\phistar)=\frac{1}{2\pi}\sum_n c_n e^{in\phistar}$ is restricted
by $\Zthree$-invariance to $c_n=0$ unless $3|n$, giving
$W(\phistar)\propto 1+a_3\cos(3\phistar)+a_6\cos(6\phistar)+\cdots$.
The leading non-zero coefficient $a_3\neq0$ reflects the 120°
crossing geometry; higher harmonics $a_6,a_9,\ldots$ are
suppressed by powers of $r_0/R_0\ll1$.
\qed
\end{proof}

\begin{remark}[Contrast with photon-initiated production]
\label{P18:rem:photon_contrast}
For $e^+e^-\to\gamma^*\to\rho^+\rho^-$ (photon-initiated, no
$\QH=3$ parent), the two $\rho$ mesons are produced from the
photon's $J^{PC}=1^{--}$ quantum numbers.  The $\phistar$
distribution is then governed by transverse photon helicity:
\begin{equation}
  W_{\gamma}(\phistar) = \frac{1}{2\pi}\bigl[
    1 + a_2^\gamma\cos(2\phistar)
  \bigr],
  \label{P18:eq:phi_star_photon}
\end{equation}
with $a_2^\gamma\neq0$ (from the transverse photon polarisation)
and $a_1=a_3=0$ (no $\Zthree$ geometry in the initial state).
The distinction between $\cos(3\phistar)$ (baryon-initiated) and
$\cos(2\phistar)$ (photon-initiated) is unambiguous and parameter-free.
\end{remark}

\subsection{Coplanarity at symmetric three-jet events}
\label{P18:sec:coplanarity}

A complementary prediction concerns the three-jet topology when
all three crossings are the production vertices simultaneously
(baryon $\to$ three quarks/jets).

\begin{proposition}[Three-jet coplanarity from B-departure geometry]
\label{P18:prop:coplanarity}
The three primary jets from a $\QH=3$ baryon exit at the three
B-segment departure zones, immediately past the three crossing
vertices.  By Theorem~\ref{P18:thm:crossing_geometry}(i), the tangent
at each crossing lies exactly in the $z=+r_0$ horizontal plane
(zero $z$-component), and the three departure directions are
separated by $2\pi/3$.  The three primary jets therefore lie in
a common elevated plane $z=+r_0$, giving zero intrinsic
acoplanarity at the topological level.  Residual acoplanarity in
LEP three-jet events~\cite{Abe1997,Acton1993,Botner1992} is
attributed to perturbative soft-gluon radiation, consistent with
this prediction.
\end{proposition}

%══════════════════════════════════════════════════════════════════════════════
\section{Connection to the Arc-Segment Energy Partition}
\label{P18:sec:arcseg}
%══════════════════════════════════════════════════════════════════════════════

Paper~XVI~\cite{Paper16} computes the $J_4$ energy partition of
$T_{2,3}$ into six arc-segments (three A-segments of 90° each,
three B-segments of 30° each) and records the three categories
of field disturbance in the 12-site fraying pattern.  We connect
these numerical findings to the topological framework above.

\begin{proposition}[A-segment dominance]
\label{P18:prop:arcseg_dominance}
The A-segments (midpoint-to-crossing, 90°) carry $\approx79\%$
of the total $J_4$ charge at energy density $\rho_A/\rho_B\approx1.23$
(Paper~XVI~\cite{Paper16}, Table~1, $N=96$ grid).
The Hopf flux is concentrated in the propagation arms between midpoints
and crossings, not at the crossing vertices themselves.
\end{proposition}

\begin{proposition}[$\rho_{J_4}$ anti-correlates with $S_{\mathrm{eff}}$]
\label{P18:prop:anticorrelation}
The effective suppression $S_{\rm eff}(\theta,\rho)=
\sin^4\theta/[\vp^6(1+\beta\rho)]$ is maximised at the crossing
vertices ($\kappa=0.399$, curvature maximum) and minimised at the
midpoints ($\kappa=0.026$, factor~15 below crossings).
The field $\rho_{J_4}$ is depleted precisely where $S_{\rm eff}$
is largest: the crossing vertices are topological bottlenecks
where the field is constrained to its minimum.
This anti-correlation is the condensate-level signature of the
confinement condition: where confinement is strongest
(pure colour, simple-current, $\dimq{3/2}=1$), topological flux
is minimal.
\end{proposition}

%──────────────────────────────────────────────────────────────────────────────
\subsection{Crossing vertices as topological lenses: the energy cascade}
\label{P18:sec:cascade}
%──────────────────────────────────────────────────────────────────────────────

The anti-correlation of Proposition~\ref{P18:prop:anticorrelation}
reveals that crossing vertices are not energy sources but
\emph{topological lenses}: they redirect incoming energy
from the A-segment into the B-segment departure channel.
The following proposition makes this precise.

\begin{proposition}[Energy cascade through the $T_{2,3}$ tube]
\label{P18:prop:cascade_p18}
The physical energy cascade in the $\QH=3$ condensate, as observed in
the $N=256$, $h=0.0375$ gradient-flow run of Paper~XVI~\cite{Paper16}
(Remark~12.3%\ref{P18:rem:gf_twophase_p18}
), follows the causal chain:
\begin{enumerate}[noitemsep,label=(\roman*)]
\item \textbf{Initial fraying at distal lobes.}
  The three outer distal lobe sites ($R_{xy}=R_0+r_0=3.874$, $z=0$,
  $\kappa=0.349$) are the first to show visible energy depletion,
  appearing at the outward (vacuum-facing) face of the tube as early
  as step~50 after Phase~1 exit.  The lobes sit at maximum radial
  extent, where the condensate density is most thinly distributed
  over the largest tube circumference and the gradient toward vacuum is
  most directly aligned with the outward radial direction.  This makes
  radial exposure, not curvature, the operative weakness at the lobe.
\item \textbf{Crossing depletion propagates inward.}
  Simultaneously with lobe fraying, the crossing vertices
  ($R_{xy}=R_0$, $z=+r_0$, $\kappa=0.399$) show depletion on their
  outer (vacuum-facing) side of the tube.  By step~1000 this depletion
  has propagated to the inner side.  The crossing acts as a
  \emph{topological lens}: it cannot radiate outward (the tangent lies
  exactly in the $z=+r_0$ horizontal plane), so the depletion is inward
  propagation of the structural failure arriving from the lobe side.
  No densification occurs at the crossings at any stage.
\item \textbf{Redirection at B-segment departure (primary jets 1, 2, 3).}
  Energy that does not exit at the lobe is redirected at the crossing
  into the B-segment departure zone, propagating radially outward in
  the $z=+r_0$ horizontal plane, 120° from its two partners.
  The three primary jets are coplanar in the $z=+r_0$ plane.
\item \textbf{Secondary lobe emission (jets 4, 5, 6).}
  Energy propagating along the B-segment arc (30°) reaches the distal
  lobe at $R_{xy}=R_0+r_0$, $z=0$, with tangent z-component
  $+0.321|\dot\Gamma|$ (upward from $z=0$, below the three-jet plane
  at $z=+r_0$).  Secondary emission here is deflected
  $\approx16^\circ$--$19^\circ$ from the three-jet event plane
  (Section~\ref{P18:sec:fourjet}).
\item \textbf{Midpoints last.}
  The three midpoint torsion-twist sites ($R_{xy}=R_0-r_0=2.126$,
  $z=0$, $\kappa=0.026$) show only minor perturbation throughout,
  remaining the last and least disturbed category.  Their inner
  radial position ($R_{xy}$ minimum) provides the most protection
  from vacuum exposure; their $S_{\rm eff}$ minimum reflects
  geometric weakness, but this is secondary to radial exposure in
  determining the empirical onset order.
\end{enumerate}
The crossing vertex is an energy \emph{redirector}, not an energy
\emph{source}.  The three conventional quark jets exit from the
B-segment departure zones.  The distal lobes are the primary
structural failure sites; the midpoints are the most protected.
\end{proposition}

\begin{remark}[Implosion, then explosion]
\label{P18:rem:implode_explode}
The $N=256$ snapshots of Figure~\ref{P18:fig:gf_n256_cascade} support a
two-stage decay picture sharper than ``lobes fray first.''  At
step~$1000$ the three lobes deplete together but diffusely, with no
organised escape route; by step~$2000$--$3000$ the depleted energy is
\emph{not} found in the vacuum outside the tube but reflected inward
along the A-segment (the inward-deflecting behaviour already noted
following equation~\eqref{P18:eq:lobe_angle} and in the discussion of the
$N=256$ run above).  This is an \emph{implosion}: the three lobes,
sharing the same maximal radial exposure, collapse together under the
surrounding density gradient, and the resulting energy is funnelled
inward rather than escaping outward, because the topological floor
forbids the field from simply unwinding into vacuum.

By step~$8000$ (Figure~\ref{P18:fig:gf_n256_8000}) the field has
reorganised into sharply bounded bright arcs against clearly emptied
gaps --- the signature of an \emph{explosion} phase, in which the
energy displaced inward by the lobe implosion is redirected through
the crossings.  The crossings are the loci of \emph{maximum}
$S_{\rm eff}$ (Proposition~\ref{P18:prop:anticorrelation}): the
density-feedback suppression is strongest there, and $\rho_{J_4}$
is consequently depleted to its minimum.  The crossings are not
release valves where suppression is weak; they are topological
bottlenecks where suppression is strongest, through which energy
is forced because the topological floor forbids any direct
radial escape to vacuum.  The crossing acts as a
\emph{redirector} throughout: during the implosion phase
($n\lesssim3000$) it receives energy funnelled inward from the
lobes; during the explosion phase ($n\gtrsim3000$) it pushes that
energy into the B-segment departure channels, where the field can
finally exit.  The crossings are thus never sources in the
energetic sense --- they are conduits, and the energy that
emerges at the B-departure zones was originally carried by the
lobes.

The trefoil vortex knots realised in water by Kleckner and
Irvine~\cite{Kleckner2013} offer a partial, but only partial,
analogy.  There, reconnection is observed where two strands of the
knot draw close, become nearly parallel, and twist around one another
before colliding~\cite{Kleckner2013} --- a description that maps most
naturally onto the \emph{crossings} of $T_{2,3}$, where the two
strands of the condensate trefoil pass nearest each other, not onto
the lobes.  The $\QH=3$ condensate's depletion instead begins at the
lobes, the points of maximal radial exposure to vacuum, and only later
reaches the crossings.  This is a genuine point of disanalogy, not a
detail to be smoothed over: the two systems share a qualitative
two-stage character (a build-up phase followed by a localised release)
but differ in which geometric feature of the knot triggers the first
stage --- proximity of strands in the fluid case, radial exposure to
vacuum in the condensate case, a difference traceable to the
density-feedback term $\sin^4(f)$ having no counterpart in ordinary
viscous reconnection (Remark~\ref{P18:rem:saddle_difficulty}).
\end{remark}

\begin{remark}[Cascade ordering]
\label{P18:rem:cascade_ordering}

The $N=256$ gradient-flow data shows directly that the operative
weakness is \emph{radial exposure} (distance from the tube axis
relative to the vacuum boundary), not curvature.

The midpoints at $R_{xy}=R_0-r_0=2.126$ are geometrically unconstrained
(low $\kappa$) but radially protected.  The distal lobes at
$R_{xy}=R_0+r_0=3.874$ are geometrically intermediate ($\kappa=0.349$)
but radially exposed.  Under gradient flow that minimises $E_{\rm geom}$
globally, the field finds it cheapest to shed energy at the most exposed
sites first.

Conjecture~\ref{P18:conj:six_quarks} (six quarks from midpoint-breaking
combinatorics) The midpoints remain the structurally weakest sites under
Frenet geometry and the most likely breaking sites under thermal or quantum
fluctuations beyond gradient flow. The gradient-flow ordering reflects
the specific path of steepest descent from the family boundary, not the
equilibrium stability hierarchy.
\end{remark}

\begin{remark}[Why the saddle point is elusive in Paper~XVI]
\label{P18:rem:saddle_difficulty}
Proposition~\ref{P18:prop:anticorrelation} and
Theorem~\ref{P18:thm:trefoil_satellite_jones} together explain a structural
difficulty reported in Paper~XVI~\cite{Paper16}: the gradient-flow
algorithm converges to the $\QH=3$ energy-density field but does
not locate a unique sharp BPS saddle.

At $\QH=2$, the BPS saddle condition $K_{\rm fb}/J_4=\vp^6/2^{1/3}$
has a unique solution $C^*$ because $\vp$ is irrational: the
condition pins the profile to a single algebraically irrational
fixed point.  The irrationality is the anchor.

At $\QH=3$, the satellite Jones invariant of the trefoil
(Theorem~\ref{P18:thm:trefoil_satellite_jones}) satisfies
$|\Jones_{\rm sat}(q_5)|=1/\vp$ (irrational --- the same modulus as
the $\QH=2$ Hopf link's own invariant, related to it by a single
phase, Remark~\ref{P18:rem:richer_criterion}), together with the
simple-current condition $\dimq{3/2}=1$.  Since the trefoil's own
invariant is irrational, the argument used at $\QH=2$ does not, on
its own, explain why gradient flow fails to locate a sharp $\QH=3$
saddle.  One conjectural candidate explanation is that the physical
\emph{formation path} to the trefoil passes through an unfused
intermediate configuration (the $n=4$ member of the satellite twist
family, Proposition~\ref{P18:prop:satellite_twist_family}), which
does have a fully rational satellite invariant
$\Jones_{\rm sat}^{(4)}(q_5)=1$ (equation~\eqref{P18:eq:sat_n4}); if
so, it would be the rationality of this intermediate, not the
irrational value of the completed trefoil, that removes the
irrational algebraic anchor during formation.  This is \emph{not}
established: nothing here demonstrates that the twist parameter $n$
corresponds to a physically continuous coordinate of the actual
formation dynamics, as opposed to being a convenient parametrisation
of a family of diagrams.

The gradient flow nonetheless converges to the $\QH=3$ family
(explaining why the arc-segment partition, fraying pattern, and
$\rho_{J_4}$--$S_{\rm eff}$ anti-correlation are all reproducible
and stable) without converging to a specific point within it
(explaining why a unique $C^*$ is not found).  Whether this
elusiveness reflects the rational $n=4$ intermediate conjectured
above, some other feature of the formation dynamics, or is better
understood some other way remains open.

The gradient flow instead settles at the
family boundary $K_{\min}\approx2289$ (Paper~XVI~\cite{Paper16},
Proposition~\ref{P18:prop:anticorrelation}), which is a genuine physical
quantity --- the lowest energy accessible within the flat family ---
rather than an irrational BPS fixed point.

A consequence of the flatness is that the minor radius $r_0=0.874$
is set by the profile construction ansatz rather than by a BPS
minimisation.  Paper~XVII~\cite{Paper17} (Remark~11.2 %\ref{P17:rem:solar_angle}
)
identifies the sharp condition $r_0=2R_0/(3\vp+2)=0.8754$ under
which the midpoint emission angle of $T_{2,3}$ equals the leading-order
solar mixing angle $\arctan(1/\vp)$ exactly ($0.16\%$ match).  A
tighter candidate, $r_0=R_0/C_2^*=0.8742$ ($0.02\%$ match), is
motivated by the satellite construction
(Proposition~\ref{P18:prop:satellite},
Remark~\ref{P18:rem:satellite_physical}): if $1/C_2^*$ is the
dimensionless tube radius of $N(\gamma_1)$, $r_0$ inherits its scale
from the $\QH=2$ profile-sharpness constant $C_2^*=3.4318$ derived
in Paper~I~\cite{Paper1}, Theorem~3.1 and Paper~XV~\cite{Paper15},
Remark~6.1.  %\ref{P15:rem:r0_origin}
Neither candidate is derived from the $\QH=3$ Euler--Lagrange
equations; the flat landscape admits both.

The two forms of the solar-angle candidate are related by the
standard Fibonacci identity $\vp^4=3\vp+2$, giving
$2R_0/(3\vp+2) = 2R_0/\vp^4 = 6/\vp^4$ at $R_0=3$.
Substituting the proved Bogomolny parameter $\lambda_3=\vp^6$
yields the equivalent form
\begin{equation}
  \lambda_3\cdot r_0 \;=\; 2R_0\cdot\vp^2,
  \label{P18:eq:lambda3_r0_check}
\end{equation}
which follows algebraically:
$\vp^6\cdot(2R_0/\vp^4)=2R_0\vp^2$.  This shows that the solar-angle
candidate $r_0=2R_0/(3\vp+2)$ is self-consistent with the
independently-derived $\lambda_3=\vp^6$: two constants fixed by
different arguments of the framework satisfy~\eqref{P18:eq:lambda3_r0_check}
identically.  With the numerical value $r_0=0.874$ (rather than the
exact candidate $0.8754$), the left side of~\eqref{P18:eq:lambda3_r0_check}
differs from the right by $0.16\%$, the same discrepancy as the
$r_0$ comparison itself --- as expected, since the two statements are
algebraically equivalent.

The flat-bottomed landscape (an established numerical fact, Papers~XV/XVI
gradient-flow searches; the rational-$n=4$-intermediate account of
\emph{why} it is flat, above, remains conjectural) means that
gradient-flow minimisation cannot select a preferred $r_0$: the
energy is monotone in $r_0$ over the candidate range, with no minimum.  Equation~\eqref{P18:eq:lambda3_r0_check} is
therefore a \emph{consistency} statement, not a falsifiability
criterion: it shows the constants do not contradict each other, but
the $\QH=3$ gradient flow cannot decide between the two candidates.
The decisive test would be a first-principles derivation of $r_0$
from the $\QH=2$ tube geometry via the satellite construction
(Proposition~\ref{P18:prop:satellite}), not a numerical minimisation.

\textbf{This instability is not a numerical artefact of the kind that
afflicts knotted vortices in a real fluid.} Laboratory trefoil vortex
knots, realised by Kleckner and Irvine~\cite{Kleckner2013} (2013) in water, unknot within a
few hundred milliseconds via a localised reconnection event driven by
viscosity --- an imperfection of the medium that an ideal, inviscid
fluid would not have. The $\QH=3$ condensate's instability has the
opposite origin: it is a property of the medium's own energy
functional, not a departure from an ideal one. The density-feedback
term $\sin^4(f)$ is what gives the $\QH=1$ sector its topological
floor, the energy floor below which no further unwinding is possible,
while leaving $\QH=3$ free to relax toward the $(2{+}1)$ configuration
identified in Section~\ref{P18:sec:z3_drift}, in the same qualitative
direction the trefoil vortex took in the tank. The geometry, including
the feedback term, prefers the lighter configuration; that preference,
not a flaw in the medium, is what fixes where the $\QH=3$ sector
settles.
\end{remark}

\subsection{Gradient-flow cascade: $N=192$ symmetry breaking and $N=256$ cascade ordering}
\label{P18:sec:z3_drift}
Two gradient-flow runs reported in Paper~XVI~\cite{Paper16}
(Remark~12.3%\ref{P18:rem:gf_twophase_p18}
) provide complementary evidence for the structure of the $\QH=3$ energy
landscape.

\textbf{$N=192$ run: spontaneous $\Zthree$ breaking.}
At step $n=4000$ ($J_4\approx1600$) the field is approximately
$\Zthree$-symmetric.  At step $n=6000$ ($J_4\approx702$) the symmetry
is spontaneously broken: two outer lobe sites carry elevated $\rho_{J_4}$
while the third outer lobe and its adjacent midpoint are depleted.
Energy has concentrated on either side of the depleted region.
This $(2{+}1)$ lobe splitting arises from floating-point accumulation in
Adam's moment estimates ($\sim10^{-7}$ per step) on the flat valley floor
(Remark~\ref{P18:rem:saddle_difficulty}); the particular direction selected is
seed-dependent.  The observation demonstrates that the flat landscape
admits $(2{+}1)$ deformations at near-zero energy cost, consistent with
Conjecture~\ref{P18:conj:six_quarks}.

\begin{figure}[htp]
    \centering
        \begin{subfigure}[T]{0.45\linewidth}
        \centering
        \includegraphics[width=\linewidth]{n_2000_density.png}
        \caption{Step $n=2000$ (Phase~1 exit / early Phase~2):
        energy-density view, $J_4=2834$, directional arrows
        off, density threshold $21\%$.
        The $\rho_{J_4}$ distribution is approximately
        $\mathbb{Z}_3$-symmetric: all three outer lobe sites carry
        comparable flux, and the three A-segment tube bodies
        are approximately uniformly bright.  The initial breaking can
        be observed at the midpoint on the right arc. And 2 outer lobe
        regions on the left upper and lower arcs can be seen concentrating
        energy, consistent with the
        $\rho_{J_4}$--$S_{\rm eff}$ anti-correlation
        (Proposition~\ref{P18:prop:anticorrelation}).
        This configuration represents the field near the
        $K$-minimum ($K_{\min}\approx2289$), the boundary of the
        flat-bottomed $\QH=3$ energy family
        (Remark~\ref{P18:rem:gf_twophase_p18}).}
        \label{P18:fig:gf_n192_4000}
    \end{subfigure}
    \hfill
    \begin{subfigure}[T]{0.45\linewidth}
        \centering
        \includegraphics[width=\linewidth]{n_6000_arrows.png}
        \caption{Step $n=6000$ (Phase~2 drift): directional-arrow
        view, $J_4=702$, density threshold $6\%$.
        Arrow direction encodes $\mathbf{n}(\mathbf{x})$;
        arrow color encodes $\rho_{J_4}$.
        The $\mathbb{Z}_3$ symmetry of Construction~11.27 of Paper~XVI~\cite{Paper16}
        is demonstrably broken: two outer lobe sites (upper-left and
        lower arcs) retain dense, coherent flux with high density arrows,
        while the third outer lobe and its adjacent midpoint (right
        arc, approximately $3$ o'clock) are depleted, with
        less coherent and dense arrows.  Energy has concentrated
        on either side of the depleted regions; the crossing
        sites show asymmetric $\rho_{J_4}$ depletion around
        the tube.  This $(2{+}1)$ lobe splitting is the spatial
        signature of a single-midpoint string-breaking mode
        (Section~\ref{P18:sec:z3_drift},~\ref{P18:rem:gf_twophase_p18}).}
        \label{P18:fig:gf_n192_6000}
    \end{subfigure}
    \caption{$\QH=3$ condensate field at steps $n=2000$ and
    $n=6000$ of the \textbf{$N=192$, $h=0.05$ run} (the lower-resolution
    run in which Adam-optimiser floating-point noise drives a
    spontaneous, asymmetric $\Ztwo$-breaking event), visualised by
    \texttt{hopfion\_viewer.html} using $1.5\times10^5$ particles
    ($R_0=3$, $r_0=0.874$, $C^*=2.5062$).
    \emph{Left}: energy-density mode; particle brightness
    encodes $\rho_{J_4}$, hue encodes $\mathbf{n}$ direction.
    \emph{Right}: directional-arrow mode; each arrow represents
    a sample point with direction $\mathbf{n}$ and color
    proportional to $\rho_{J_4}^{1/2}$.
    The transition from approximate $\mathbb{Z}_3$ symmetry (left)
    to spontaneously broken $(2{+}1)$ lobe distribution (right)
    occurs during Phase~2 on the flat-bottomed $\QH=3$ energy
    landscape (Remark~12.3 of Paper~XVI~\cite{Paper16}): %\ref{P16:rem:gf_twophase}
    the gradient flow drifts from the $\mathbb{Z}_3$-symmetric locus
    along a $(2{+}1)$ flat direction corresponding to
    single-midpoint string breaking.  One conjectural (not established,
    Remark~\ref{P18:rem:saddle_difficulty}) account of the absent
    restoring force is that the formation event passes through the
    unfused intermediate, which has rational satellite invariant
    $\Jones_{\rm sat}^{(4)}(q_5)=1$
    (Proposition~\ref{P18:prop:satellite_twist_family}), leaving no
    irrational algebraic anchor during formation.  This $N=192$ run is retained
    here as a record of the asymmetric $(2{+}1)$ breaking mode
    actually observed at lower resolution: it is a genuine alternative
    deformation accessible on the flat valley floor
    (Remark~\ref{P18:rem:saddle_difficulty}), distinct from the
    symmetric three-lobe collapse seen at higher resolution
    (Figure~\ref{P18:fig:gf_n256_cascade}), and is consistent with
    Conjecture~\ref{P18:conj:six_quarks}'s combinatorics of
    single-site versus multi-site breaking patterns.}
    \label{P18:fig:gf_n192_snapshots}
\end{figure}

\textbf{$N=256$ run: direct observation of fraying site ordering.}
The higher-resolution run ($N=256$, $h=0.0375$, $\delta_{\max}=2^\circ$)
suppresses the Adam-induced symmetry breaking and reveals the fraying
onset order directly (Figure~\ref{P18:fig:gf_n256_cascade}).  At step~$50$
the field is dense and approximately $\Ztwo$-symmetric, immediately
following Phase~1 exit.  By step~$1000$, $J_4$ has fallen $42\%$ and
depletion is diffuse, with the three outer lobes shedding together
rather than any single category of site standing out; crossing sites
show simultaneous outer-tube depletion, while midpoint disturbance
remains negligible.  By step~$3000$, the depletion has organised into
a continuous pattern rather than a diffuse one, and the energy density
at the lobes is densifying on both sides of the break while deflecting
\emph{inward} rather than outward --- the onset of the implosion phase
of Remark~\ref{P18:rem:implode_explode}.  By step~$8000$, the field has
reorganised into sharply bounded bright arcs against emptied gaps with
large local reorientation of $\mathbf{n}$, consistent with the
explosion phase in which energy displaced inward by the lobe collapse
is redirected through the crossing bottlenecks (maximum-$S_{\rm eff}$
sites, Proposition~\ref{P18:prop:anticorrelation}) and exits at the
B-departure zones.  Midpoints remain
the least disturbed category throughout the run.

The empirical ordering (lobes deplete first and together; crossings
follow, first as redirectors then as the explosion-phase source;
midpoints last) is incorporated into the revised
Proposition~\ref{P18:prop:cascade_p18} and the implosion--explosion picture of
Remark~\ref{P18:rem:implode_explode}.  The observation that no energy
exits outward during the implosion phase is direct numerical evidence
for the topological barrier of the $\QH=3$ sector: the field reaches
the boundary of the flat family and reflects inward, confirming that
confinement operates even within the gradient-flow dynamics, before
the delayed outward release at the crossings during the explosion
phase.

\begin{figure}[htp]
    \centering
    \begin{subfigure}[T]{0.23\linewidth}
        \centering
        \includegraphics[width=\linewidth]{n_000050.png}
        \caption{Step $n=50$: $J_4=5602$, $\bar r=3.595$.
        The tube is dense and approximately $\mathbb{Z}_3$-symmetric;
        all three lobes, all three crossings, and the A-segment
        bodies appear comparably bright.  This is the field
        immediately following Phase~1 exit, before any visible
        depletion.}
        \label{P18:fig:gf_n256_50}
    \end{subfigure}
    \hfill
    \begin{subfigure}[T]{0.23\linewidth}
        \centering
        \includegraphics[width=\linewidth]{n_001000.png}
        \caption{Step $n=1000$: $J_4=3239$, $\bar r=3.570$.
        $J_4$ has fallen by $42\%$.  Depletion is diffuse and
        distributed around the full loop rather than concentrated
        at a single category of site; no single lobe, crossing, or
        midpoint yet stands out as uniquely emptied.  This is
        consistent with the radial-exposure mechanism of
        Proposition~\ref{P18:prop:cascade_p18}: all three lobes, sharing the
        same maximal radial extent, begin shedding together.}
        \label{P18:fig:gf_n256_1000}
    \end{subfigure}
    \hfill
    \begin{subfigure}[T]{0.23\linewidth}
        \centering
        \includegraphics[width=\linewidth]{n_003000.png}
        \caption{Step $n=3000$: $J_4=2717$, $\bar r=3.566$.
        A continuous bright arc has organised across roughly half
        the structure, with the remainder visibly sparser.  The
        depletion pattern has localised: this is the onset of the
        \emph{implosion} phase, in which the tube collapses inward
        from the lobes rather than radiating outward into vacuum
        (Remark~\ref{P18:rem:implode_explode}).}
        \label{P18:fig:gf_n256_3000}
    \end{subfigure}
    \hfill
    \begin{subfigure}[T]{0.23\linewidth}
        \centering
        \includegraphics[width=\linewidth]{n_008000.png}
        \caption{Step $n=8000$: $J_4=1406$, $\bar r=3.561$.
        $J_4$ has halved again.  The surviving flux is organised
        into sharply-bounded bright arcs against clearly emptied
        gaps, with strong angular excursions in $\mathbf{n}$
        (colour) indicating large local reorientation rather than
        uniform fading.  This is consistent with the \emph{explosion}
        phase: energy displaced inward by the lobe implosion is
        redirected through the crossing bottlenecks
        (maximum-$S_{\rm eff}$ sites,
        Proposition~\ref{P18:prop:anticorrelation}) and exits at the
        B-departure zones
        (Remark~\ref{P18:rem:implode_explode}).}
        \label{P18:fig:gf_n256_8000}
    \end{subfigure}
    \caption{$\QH=3$ condensate field at steps $n=50$, $1000$,
    $3000$, $8000$ of the $N=256$, $h=0.0375$ gradient-flow run,
    visualised by \texttt{hopfion\_viewer.html} using
    $1.5\times10^5$ particles ($R_0=3$, $r_0=0.874$, $C^*=2.5062$).
    Energy-density mode throughout; particle brightness encodes
    $\rho_{J_4}$, hue encodes $\mathbf{n}$ direction.  Unlike the
    $N=192$ run of Figure~\ref{P18:fig:gf_n192_snapshots}, the higher
    resolution suppresses Adam-induced floating-point symmetry
    breaking and shows the three outer lobes depleting together,
    not asymmetrically: this is the direct numerical evidence for
    Proposition~\ref{P18:prop:cascade_p18}'s revised cascade ordering
    (lobes first, simultaneously) and for the
    implosion--explosion picture of
    Remark~\ref{P18:rem:implode_explode}.}
    \label{P18:fig:gf_n256_cascade}
\end{figure}

\begin{remark}[Three remarks on interpretation.  First, the solver computes]
\label{P18:rem:three_remark_interpretation}
$E_{\rm geom}=K\!\cdot\!J_4$ as a global sum; the redistribution is
genuine global dynamics, not a per-segment artefact.  Second, the
revised ordering does not invalidate Conjecture~\ref{P18:conj:six_quarks}:
the midpoints remain the geometrically weakest sites under Frenet
analysis, and under thermal or quantum fluctuations they are still the
expected breaking sites; the gradient-flow ordering reflects the steepest
descent from the family boundary, not the equilibrium stability hierarchy.
Third, the near-perfect $\Zthree$ symmetry of the $N=256$ run throughout
Phase~2 confirms that the $(2{+}1)$ breaking in the $N=192$ run was a
numerical artefact of the larger per-step rotations.
\end{remark}


\begin{remark}[Two-phase landscape, Z3-symmetry breaking, and the Jones
  polynomial explanation]
\label{P18:rem:gf_twophase_p18}
The data from the $N=216$, $h=0.05$  with:
Phase 2: $lr=5^{-5}$, $\Delta{max}=5.0°$ per point and
Phase transition: K rises > 1.0 above its minimum,
run confirm the two-phase structure
and refine the sector boundary estimate to $K_{\min}\approx2289$
(improved from $2294$, $N=192$, $h=0.05$ with:
Phase 2: lr=0.0001, $\Delta{max}=9.0°$ per point
Phase transition: K rises > 2.0 above its minimum), consistent with resolution
convergence toward the continuum value).

\emph{Phase~1} genuinely operates within the $\QH=3$ sector: the field
sharpens without topology change, $\bar r$ drifts toward $R_0=3$,
and the gradient norm approaches a natural plateau.  The $K$-minimum
marks the closest numerical approach to the $\QH=3$ sector boundary
before the field begins to drift toward the vacuum.

\emph{Phase~2} exposes a structural feature of the $\QH=3$ landscape that
is not numerical in origin.  Visual inspection of the energy-density field
$\rho_{J_4}$ at steps $n=4000$ and $n=6000$ (Figure~\ref{P18:fig:gf_n192_snapshots})
reveals that the field spontaneously breaks the exact $\mathbb{Z}_3$ symmetry
of Paper~XVI~\cite{Paper16}, Construction~11.27 during Phase~2.  At step
$n=4000$ the distribution is approximately $\mathbb{Z}_3$-symmetric.  By
step $n=6000$, with $J_4$ halved to $702$, the symmetry is broken: two of
the three outer lobe sites carry elevated $\rho_{J_4}$, while the third
outer lobe and its adjacent midpoint have lost energy.  Energy has built
up on either side of the depleted region, and the crossing sites show
asymmetric depletion consistent with the $\rho_{J_4}$--$S_{\rm eff}$
anti-correlation (Proposition~\ref{P18:prop:anticorrelation}) operating
non-uniformly around the tube.  The overall energy distribution is not
$\mathbb{Z}_3$-symmetric.

The mechanism is the following.  Construction~11.27 of Paper~XVI~\cite{Paper16}
is exactly $\mathbb{Z}_3$-symmetric, and Phase~1 preserves this because
the gradient of $E_{\rm geom}$ at a $\mathbb{Z}_3$-symmetric configuration
is itself $\mathbb{Z}_3$-symmetric.  Phase~2, however, operates on the
flat-bottomed energy valley discussed in
Remark~\ref{P18:rem:saddle_difficulty}: without an irrational
algebraic anchor to pin the profile to a unique minimum, the
landscape is locally flat and no restoring force returns the field
to the symmetric locus after any perturbation.  Floating-point
accumulation in Adam's adaptive moment estimates breaks $\mathbb{Z}_3$
at the $\sim10^{-7}$ level over thousands of steps; once off the symmetric
locus, the gradient flow drifts along the asymmetric flat directions with
no restoring force.

The observed asymmetry pattern---two lobes retaining energy, one lobe and
its adjacent midpoint losing energy---is a $(2{+}1)$ splitting of the
trefoil's three outer lobes.  This is precisely the pattern predicted by
Conjecture~\ref{P18:conj:six_quarks} for the lightest single-midpoint
string-breaking mode: breaking one of the three midpoint torsion-twist sites
produces a $(2{+}1)$ separation of the tube's $\mathbb{Z}_3$ structure.  The
gradient flow is spontaneously selecting one of the three symmetry-equivalent
$(2{+}1)$ directions on the flat valley floor.  This is not a coincidence in
the sense that the Jones polynomial rationality forces the field onto a
flat landscape where all three $(2{+}1)$ directions are accessible; it is
physically significant that the drift is toward a well-defined string-breaking
mode rather than an arbitrary deformation.
-
Grid refinement addresses the numerical resolution problem; the structural
$\mathbb{Z}_3$-symmetry breaking during Phase~2 is a topological feature of
the $\QH=3$ flat landscape and will persist at any resolution.
\end{remark}

\begin{remark}[$K_{\min}$ resolution convergence and three-point extrapolation]
\label{P18:rem:kmin_extrapolation}
The $K_{\min}$ values from three gradient-flow runs at increasing resolution
are:
\begin{center}
\begin{tabular}{cccc}
\toprule
$N$ & $h$ & $K_{\min}$ & Points across tube \\
\midrule
144 & 0.075 & 2313.9 & 10.6 \\
192 & 0.050 & 2289\phantom{.0} & 16.0 \\
256 & 0.0375 & 2268\phantom{.0} & 21.3 \\
\bottomrule
\end{tabular}
\end{center}
The successive drops are $\Delta K_1=-24.9$ ($N=144\to192$) and
$\Delta K_2=-21.0$ ($N=192\to256$), with $h$-ratios $h_1/h_2=1.5$ and
$h_2/h_3=4/3$.  Fitting the power law
$K_{\min}(h)\approx K_{\min}^\infty+A\,h^\alpha$ to the three points
gives $\alpha\approx0.75\pm0.15$ and
\begin{equation}
  K_{\min}^\infty \;\approx\; 2230\text{--}2245,
  \label{P18:eq:kmin_extrap}
\end{equation}
with the uncertainty reflecting the range of $\alpha$ consistent with the
data.  The convergence is monotone from above, consistent with an
$O(h^\alpha)$ discretisation error with $\alpha$ slightly below~1 (likely
from the interplay of second-order finite differences and the
Construction~C initial condition's residual grid-scale structure).
The continuum $K_{\min}^\infty$ is the family-boundary energy of the
$\QH=3$ sector in grid units; its conversion to physical units requires
$L_{\rm cond}$ (open problem of Paper~XVI~\cite{Paper16},
Section~16.1 %\ref{P16:sec:saddle_prerequisite}).
\end{remark}

%──────────────────────────────────────────────────────────────────────────────
\subsection{The Botner measurement framework and $\beta_3$}
\label{P18:sec:botner}
%──────────────────────────────────────────────────────────────────────────────

The DELPHI measurement of three-jet event orientation in $Z^0$
hadronic decays~\cite{Botner1992} establishes the mathematical
framework within which the topological $\phistar$ prediction
should be tested.  The differential cross-section for
$e^+e^-\to\gamma/Z\to q\bar{q}g$ in the helicity frame is:
\begin{equation}
  \frac{16\pi}{3}\frac{d^2\sigma}{d\cos\theta\,d\chi}
  = (1+\cos^2\!\theta)\sigma_U + 2\sigma_L\sin^2\!\theta
    + 2\sigma_T\sin^2\!\theta\cos 2\chi
    - 2\sqrt{2}\,\sigma_I\sin 2\theta\cos\chi,
  \label{P18:eq:botner_full}
\end{equation}
where $\theta$ is the polar angle between the electron beam and
the thrust axis, and $\chi$ is the azimuthal angle between the
three-jet event plane and the plane of the thrust axis and beam.
Integrating over $\cos\theta$ gives:
\begin{equation}
  \frac{1}{\sigma}\frac{d\sigma}{d\chi}
  = \frac{1+\beta\cos 2\chi}{2\pi},
  \qquad
  \beta = \frac{\sigma_T/\sigma_U}{1+\sigma_L/\sigma_U},
  \label{P18:eq:botner_chi}
\end{equation}
a $\cos 2\chi$ modulation from the $\mathbb{Z}_2$ symmetry of
the spin-1 photon/Z production mechanism.

\begin{proposition}[Topological $\phistar$ distribution in Botner form]
\label{P18:prop:botner_analog}
By direct analogy with~\eqref{P18:eq:botner_chi}, the $\phistar$
distribution for baryon-initiated $\rho^+\rho^-$ production is:
\begin{equation}
  \frac{1}{\sigma}\frac{d\sigma}{d\phistar}
  = \frac{1+\beta_3\cos 3\phistar}{2\pi},
  \label{P18:eq:botner_phiS}
\end{equation}
where the coefficient $\beta_3$ is determined by the arc-segment
geometry of $T_{2,3}$ rather than by electroweak parameters.
Its natural estimate from the arc-segment energy ratio is:
\begin{equation}
  \beta_3^{\rm est} \;\sim\;
  \frac{J_4^A - J_4^B}{J_4^A + J_4^B}
  = \frac{79\% - 21\%}{79\% + 21\%}
  = \frac{58}{100} = 0.58,
  \label{P18:eq:beta3_est}
\end{equation}
where $J_4^A/J_4^B \approx 3.76$ from Paper~XVI~\cite{Paper16}.
This is a geometric prediction with no free parameters:
$\beta_3$ is fixed by the $T_{2,3}$ energy partition.

The key distinction from~\eqref{P18:eq:botner_chi}:
photon/Z-initiated production gives $\cos 2\chi$
($\mathbb{Z}_2$ from spin-1 photon), while baryon-initiated
production gives $\cos 3\phistar$ ($\Zthree$ from $T_{2,3}$
trefoil).  These are orthogonal signals with different
angular frequencies and physically distinct origins.
\end{proposition}

%──────────────────────────────────────────────────────────────────────────────
\subsection{Four-jet out-of-plane structure and the Jeremie anomaly}
\label{P18:sec:fourjet}
%──────────────────────────────────────────────────────────────────────────────

\begin{remark}[Two distinct orderings: depletion order versus jet hardness]
\label{P18:rem:two_orderings}
The implosion--explosion picture of Remark~\ref{P18:rem:implode_explode}
revises \emph{which site depletes first} in the condensate dynamics
(the lobes), but this is a different question from \emph{which
outgoing jet is hardest} in a collider final state, addressed by the
propositions below.  The two orderings need not coincide, and do not:
the lobes are the first sites to lose energy (the trigger of the
implosion), while the crossings --- the maximum-$S_{\rm eff}$
bottlenecks (Proposition~\ref{P18:prop:anticorrelation}) --- never
release energy directly but redirect it into the B-departure zones,
where the primary jets 1--3 emerge.  The crossing vertex is a
\emph{redirector} throughout (Proposition~\ref{P18:prop:cascade_p18}), and
the hardest jets exit at the B-departures, not at the crossings
themselves.  The jet-hardness ordering of
Proposition~\ref{P18:prop:jet_bound} therefore remains the relevant
ordering for collider observables (jets~1--3 hardest, at B-departure
exits coplanar at $z=+r_0$; jets~4--6 softer, at the lobes), even
though the \emph{condensate-internal} depletion sequence runs the
other way (lobes first, crossings last to show depletion because
they are topologically protected by maximum $S_{\rm eff}$, midpoints
avoided throughout due to their topological positioning in the
trefoil).  The geometric content of
Propositions~\ref{P18:prop:fourjet_oop} and~\ref{P18:prop:jet_bound} below is
unaffected by the revised cascade ordering; only its physical
narrative is sharpened by it.
\end{remark}

\begin{proposition}[Out-of-plane deflection of the 4th jet]
\label{P18:prop:fourjet_oop}
The 12-site fraying pattern of Proposition~\ref{P18:prop:jet_bound}
divides into four groups of three sites each, one group per
$\Zthree$-symmetric geometric role.  Any single group, if it acts
as an emission mechanism, produces three symmetric jets separated
by $2\pi/3$ in azimuth.  The 4th jet is therefore the \emph{first
emission from a second group}, which is a geometrically distinct
category of site.  Whether this category difference translates into
an out-of-plane deflection depends on which two groups are the
first and second to emit:

\begin{itemize}[noitemsep]
\item If jets 1--3 come from the B-segment departure exits
  ($z=+r_0$, tangent horizontal) and jets 4--6 come from any group
  at $z=0$ (lobes or midpoints), the 4th jet is out of the
  three-jet event plane in elevation.
\item If jets 1--3 come from B-departure exits and jets 4--6 come
  from the crossing approach-side entry disturbances (also
  $z=+r_0$, tangent horizontal), the 4th jet is in the same plane
  but azimuthally offset by $60°$ from jets 1--3.
\item If both groups are at $z=0$ (lobes and midpoints), the two
  groups differ in tangent elevation ($0.321|\dot\Gamma|$ at lobes
  versus $0.525|\dot\Gamma|$ at midpoints), so jets 4--6 emerge at
  a different elevation angle than jets 1--3.
\end{itemize}

Except for the (B-departure, entries) same-plane pairing, some
out-of-plane structure --- either in elevation or in emission
angle --- is generic to the 4th jet under any ordering that
alternates between distinct site groups.

The specific numerical prediction is set by the group identification.
Under the ordering of Proposition~\ref{P18:prop:jet_bound} (jets 1--3
at B-departure exits, jet 4 at the first distal lobe), the lobe
secondary emission has tangent z-component $+0.321|\dot\Gamma|$ at
$z=0$, giving:
\begin{equation}
  \Delta\theta_{\rm lobe}
  = \arcsin\!\left(\frac{r_0\cdot0.321}{R_{xy,\rm lobe}}\right)^{1/2}
  \;\approx\; 16^\circ\text{--}19^\circ,
  \label{P18:eq:lobe_angle}
\end{equation}
using $r_0=0.874$, $R_{xy,\rm lobe}=R_0+r_0=3.874$.
This is a discrete geometric prediction, not a
continuous perturbative distribution.  Standard QCD predicts
a unimodal distribution of 4th-jet angles centred near $0°$
(soft collinear emission approximately in-plane); the framework
predicts a preferred deflection near $16$--$19°$ below the
three-jet plane under the specific ordering above, and a
preferred deflection at some other discrete angle under any other
alternating ordering.  A broad, unimodal QCD-like distribution is
inconsistent with all such orderings.
\end{proposition}

\begin{remark}[The Jeremie anomaly]
\label{P18:rem:jeremie}
Jeremie et al.~\cite{Jeremie2006} studied the Bengtsson-Zerwas
(BZ) angle $\chi_{\rm BZ}$ between the plane of (jet-1,\,jet-2)
and the plane of (jet-3,\,jet-4) in four-jet events from the OPAL
detector at LEP.  They found that shape differences between
experiment and the ERT QCD matrix-element prediction concentrate
in the region $m_3+m_4>10\,\mathrm{GeV}$, where $m_3+m_4$ is
the summed invariant mass of the two least energetic jets.

The framework predicts (Proposition~\ref{P18:prop:fourjet_oop}) that
the 4th jet comes from a geometrically distinct site group from
jets 1--3, producing a discrete out-of-plane deflection at a
specific angle rather than a broad perturbative distribution.
Under the specific ordering of Proposition~\ref{P18:prop:jet_bound}
(jets 1--3 at B-departure exits, jet 4 as the first distal lobe
emission), the deflection is $\approx16$--$19°$ below the
three-jet plane; under any alternating-group ordering
(Proposition~\ref{P18:prop:fourjet_oop}), the deflection is at a
different specific angle.  A non-trivial $\chi_{\rm BZ}$ between
the plane of (jet-1, jet-2) and the plane of (jet-3, jet-4)
follows in any case where jets~1--3 and jet~4 do not share the
same site group.  Furthermore, tube-carried non-perturbative jets
are expected to have higher invariant mass than perturbative
collinear gluons, consistent with the anomaly appearing precisely
at large $m_3+m_4$.  The Jeremie anomaly is thus consistent with
the framework's generic prediction of a topologically-distinct
4th jet, whichever specific ordering of site groups is realised.
\end{remark}

\begin{proposition}[Jet multiplicity upper bound and energy ordering]
\label{P18:prop:jet_bound}
From the 12-site fraying pattern of Paper~XVI, the maximum number
of topologically distinct ejection events from a single $\QH=3$
baryon is 12.  The sites divide into four groups with distinct
geometric roles:
\begin{itemize}[noitemsep]
\item \textbf{3 sites: B-segment departure exits} (one per crossing,
  departure side).  These are the primary hard jets: the energy
  arriving from the A-segment is redirected at the crossing lens and
  exits on the B-segment departure side, propagating in the
  $z=+r_0$ horizontal plane, 120° apart.  These give jets 1, 2, 3.
\item \textbf{3 sites: distal lobe secondary emissions} (one per
  outer lobe). Energy propagating past the B-departure exit continues
  along the B-segment arc (30°) to the distal lobe, where it exits
  with tangent z-component $+0.321|\dot\Gamma|$, tilted
  $\approx16^\circ$--$19^\circ$ below the three-jet event plane.
  These give jets 4, 5, 6.
\item \textbf{3 sites: crossing approach-side entry disturbances}
  (one per crossing, approach side, 60° offset from the departure
  site at the same crossing).  The field is compressed as incoming
  A-segment energy approaches the crossing bottleneck; this entry
  disturbance is visible in the fraying pattern but represents a
  softer emission than the primary departure exit.
  These give jets 7, 8, 9.
\item \textbf{3 sites: midpoint torsion-twist residuals}
  (one per midpoint, lowest $J_4$, softest).
  These give jets 10, 11, 12.
\end{itemize}
This constitutes a topological upper bound: a single $\QH=3$
condensate soliton cannot produce more than 12 distinct jets,
in contrast to standard perturbative QCD which predicts no hard
upper limit.
The energy ordering from hardest to softest, under the assumption
that B-departure exits are the dominant primary emission channel,
is:
\begin{equation}
  \underbrace{\text{B-dep exits (3)}}_{\text{hard, coplanar at }z=+r_0}
  \;\to\;
  \underbrace{\text{Lobe (3)}}_{\approx16^\circ\text{--}19^\circ\text{ off-plane}}
  \;\to\;
  \underbrace{\text{B-dep entries (3)}}_{\text{softer, near-plane}}
  \;\to\;
  \underbrace{\text{Midpoint (3)}}_{\text{softest, residual}},
  \label{P18:eq:jet_ordering}
\end{equation}
consistent with the five-jet measurement at LEP being the
experimental ceiling (lobe, entry-disturbance, and midpoint jets
are too soft to resolve at $Z$-peak energies) and with the
prediction that above 12 jets the rate must vanish for a single
$\QH=3$ object.

Two features of this ordering are more robust than the specific
identification of each group.  First, the topological upper bound
of 12 total jets holds for any ordering; it depends only on the
site-group counting, not on which group emits hardest.  Second,
the midpoints are the softest group (jets~10--12) in any ordering:
their inner radial position ($R_{xy}=R_0-r_0=2.126$) shields them
from vacuum exposure, so they are the last sites to release energy
in the gradient-flow cascade
(Figure~\ref{P18:fig:gf_n256_cascade}) and correspondingly produce the
softest jets.  The ordering of the middle two groups (lobes vs
entry disturbances as jets~4--6 vs 7--9) depends on which emission
mechanism dominates: direct topologically-suppressed lobe release
versus entry-side compression at the crossing.  The specific
ordering above assumes lobe direct emission is more efficient
than entry-disturbance emission, consistent with the lobes'
maximum radial exposure driving early depletion in the cascade;
under any ordering, the 4th jet is out of the three-jet plane
(Proposition~\ref{P18:prop:fourjet_oop}).
\end{proposition}
\label{P18:sec:speculative}
%══════════════════════════════════════════════════════════════════════════════

The following two sections record speculative correspondences between
the $T_{2,3}$ geometry and the quark sector structure.  Both are
flagged as conjectures: they are internally consistent with the
framework and make falsifiable predictions, but lack derivations from
first principles within the density-feedback action.  The requirement
for any future derivation to reproduce these structures is recorded
as a constraint.

\subsection{Six-quark counting from midpoint-breaking combinatorics}
\label{P18:sec:quark_count}

\begin{conjecture}[Six-quark combinatorial]
\label{P18:conj:six_quarks}
The $\QH=3$ baryon decays by losing topological integrity at one
or more of its three midpoint torsion-twist sites
(Paper~XVI~\cite{Paper16}, Remark~13.4).   %\ref{P16:rem:torsion_contrast}
The six quark flavours correspond to the six distinct non-trivial
breaking patterns:
\begin{equation}
  \binom{3}{1}+\binom{3}{2} = 3+3 = 6,
  \label{P18:eq:quark_count}
\end{equation}
where $\binom{3}{k}$ counts the number of ways to break exactly
$k$ of the three midpoint sites.  Breaking all three simultaneously
returns the system to the vacuum and does not produce a quark state;
breaking zero sites is the unperturbed baryon.

The six patterns divide naturally into two isospin triplets:
the $\binom{3}{1}=3$ single-breaking patterns (one midpoint broken)
and the $\binom{3}{2}=3$ double-breaking patterns (two midpoints
broken), corresponding to the charge-$+2/3$ and charge-$-1/3$
quarks respectively, or vice versa.
\end{conjecture}

\begin{remark}[Constraints on any future derivation]
\label{P18:rem:quark_constraint}
Any derivation of the quark mass spectrum $\{m_u,m_d,m_s,m_c,m_b,m_t\}$
within this framework must satisfy:
(i)~reproduce the $3+3$ combinatorial pattern;
(ii)~assign the correct isospin ($+2/3$ vs $-1/3$) to each triplet;
(iii)~the energy cost of each breaking pattern must be derivable
from the arc-segment $J_4$ data of Paper~XVI, specifically from
the deficit $\delta J_4$ at each midpoint site relative to the
adjacent A-segment average.
\end{remark}

\subsection{Isospin from the $\Ztwo$ writhe asymmetry}
\label{P18:sec:isospin}

\begin{conjecture}[Writhe asymmetry and isospin]
\label{P18:conj:isospin}
The trefoil $T_{2,3}$ has a $\Ztwo$ writhe asymmetry: the three
crossing vertices sit at $z=+r_0$ (above the equatorial plane),
while the three midpoints and three outer-lobe distal points sit
at $z=0$ (Paper~XVI, Proposition~13.3).  %\ref{P16:prop:geom_data}
This is an exact $\Ztwo$ asymmetry: the knot is not symmetric under
$z\to-z$.

The isospin assignment $I_3=+1/2$ (charge $+2/3$) corresponds to
the crossing-vertex network ($z=+r_0$), and $I_3=-1/2$ (charge
$-1/3$) corresponds to the midpoint/distal-lobe network ($z=0$).
The physical origin is: under any external field distinguishing
the two $z$-levels, the $\Ztwo$ asymmetry selects which breaking
mode is energetically favoured, thereby selecting which quark
charge emerges.  In the absence of such a field, both isospin
values are equally probable, consistent with isospin conservation
in strong interactions.
\end{conjecture}

\begin{proposition}[Isospin-triplet mass-scale ratio from tangent geometry]
\label{P18:prop:isospin_mass_ratio}
The exit tangent at each of the three site categories has a
distinct $z$-component (Paper~XVI, Proposition~13.3): %\ref{P16:prop:geom_data}
zero at the crossing vertices ($\dot\Gamma_z=0$, exact), and nonzero
at the midpoints and distal lobes ($\dot\Gamma_z\approx-0.525$ and
$\dot\Gamma_z\approx+0.321$ respectively, in units of $|\dot\Gamma|$).
Define the dimensionless \emph{out-of-plane fraction} at a site as
$\dot\Gamma_z^2$ (with $|\dot\Gamma|=1$), so that the crossing
network exits entirely in-plane (out-of-plane fraction $0$, in-plane
fraction $1$) and the midpoint/distal-lobe network exits with
averaged out-of-plane fraction
\begin{equation}
  \langle\dot\Gamma_z^2\rangle_{M,L}
  = \tfrac12\bigl(0.525^2+0.321^2\bigr) = 0.189.
  \label{P18:eq:outofplane_avg}
\end{equation}
Under Conjecture~\ref{P18:conj:isospin}, the crossing network (charge
$+2/3$, up-type) and the midpoint/distal-lobe network (charge
$-1/3$, down-type) are proposed to set their respective mass scales
through this geometric factor:
\begin{equation}
  \frac{M_{\rm down}}{M_{\rm up}}
  \;\sim\;
  \frac{\langle\dot\Gamma_z^2\rangle_{M,L}}{1-\dot\Gamma_z^2\big|_{C}}
  = \frac{0.189}{1} = 0.189.
  \label{P18:eq:isospin_mass_ratio}
\end{equation}
Taking $M_{\rm up}$ and $M_{\rm down}$ as the geometric-mean mass
of each isospin triplet, $M_{\rm up}=(m_u m_c m_t)^{1/3}$ and
$M_{\rm down}=(m_d m_s m_b)^{1/3}$, the measured value
(PDG~\cite{PDG2024}) is
\begin{equation}
  \frac{M_{\rm down}}{M_{\rm up}}
  = \frac{(0.00467\times0.0934\times4.18)^{1/3}}
         {(0.00216\times1.27\times172.69)^{1/3}}
  = \frac{0.1222\,{\rm GeV}}{0.7795\,{\rm GeV}} = 0.157,
  \label{P18:eq:isospin_mass_ratio_obs}
\end{equation}
a $21\%$ discrepancy from the geometric estimate
\eqref{P18:eq:isospin_mass_ratio}.  This is a zero-free-parameter
estimate: both $0.525$ and $0.321$ are exact tangent
$z$-components from Paper~XVI's trefoil geometry, not fitted
quantities.  The $21\%$ residual is comparable to, and better than,
several of the mechanisms ruled out in Paper~XVI's systematic
survey (undershoots ranging from a factor of $1.7$ to $49$), and is
consistent with the always-positive residual pattern of Paper~XVI
\S12 if a further generation-universal correction (not yet
identified) applies equally to both triplets.
\end{proposition}

\begin{remark}[Scope of Proposition~\ref{P18:prop:isospin_mass_ratio}]
\label{P18:rem:isospin_mass_ratio_scope}
The identification of $M_{\rm down}/M_{\rm up}$ with the
out-of-plane/in-plane tangent-fraction ratio is \emph{motivated} by
the writhe asymmetry of Conjecture~\ref{P18:conj:isospin} but is not
\emph{derived} from the density-feedback action: no argument is
given here for why mass scale should be proportional to this
specific geometric ratio, as opposed to some other function of the
tangent data.  It is recorded because it is the first
zero-free-parameter numerical connection between the trefoil's
exact geometry and the observed two-triplet mass hierarchy, and
because the $21\%$ residual is small enough to warrant further
investigation.  It does \emph{not} address the intra-triplet
hierarchy ($m_u\ll m_c\ll m_t$ or $m_d\ll m_s\ll m_b$), which
requires either the generation-dependent string-tension
contribution of Paper~XVI \S12 or a dynamical (non-static)
treatment of the $\QH=2+\QH=1\to\QH=3$ formation event, both left
for future work.
\end{remark}

%══════════════════════════════════════════════════════════════════════════════
\section{Experimental Status and Predictions}
\label{P18:sec:experiment}
%══════════════════════════════════════════════════════════════════════════════

We collect the falsifiable predictions of this paper, ordered by
sharpness and current experimental accessibility.

\begin{enumerate}[leftmargin=2em,label=\textbf{P\arabic*.}]

\item \textbf{Process-dependent $\phistar$ in $\rho^+\rho^-$.}
The $\phistar$ distribution in $\rho^+\rho^-$ from baryon
string-breaking (e.g.\ $p\bar{p}\to\rho^+\rho^-X$ at low
$p_T$) should satisfy $(1/\sigma)d\sigma/d\phistar =
(1+\beta_3\cos 3\phistar)/2\pi$ with $\beta_3\approx0.58$
(Proposition~\ref{P18:prop:botner_analog}), peaked at
$\phistar=2\pi/3$.  The same process from $e^+e^-$ gives
$(1+\beta\cos 2\chi)/2\pi$ (Botner~\cite{Botner1992}).
The two signals have different angular frequencies (3 vs 2)
and different physical origins; they are unambiguously
distinguishable and parameter-free.

\item \textbf{Three-jet coplanarity.}
Three-jet events from $\QH=3$ baryon fragmentation are coplanar
in the $z=+r_0$ plane at the topological level
(Proposition~\ref{P18:prop:coplanarity}).  The Botner
measurement~\cite{Botner1992} of $\chi$ orientation in
$e^+e^-\to q\bar{q}g$ is consistent with this; a dedicated
comparison using events from baryon-initiated production with
NLO corrections removed is the required test.

\item \textbf{Pion emission in every baryon-formation event.}
The reaction $\QH=2+\QH=2\to\QH=3+\QH=1$ predicts pion
emission is topologically forced (Theorem~\ref{P18:thm:pion_forced}).
The ratio $\sigma(\QH=3+\QH=1)/\sigma(\QH=3)=1$ in the pure
topological channel.

\item \textbf{Baryon formation from $e^-+\pi^+\to p$ type reactions.}
The satellite construction $\QH=2+\QH=1\to\QH=3$
(Remark~\ref{P18:rem:qh2_qh1}) predicts that pion-free baryon
formation requires a $\QH=1$ initial-state object.

\item \textbf{4th-jet out-of-plane deflection.}
The 4th jet is topologically required to come from a different
site group than jets 1--3
(Proposition~\ref{P18:prop:fourjet_oop}), producing a discrete
preferred deflection angle rather than a broad perturbative
distribution.  Under the specific ordering of
Proposition~\ref{P18:prop:jet_bound} (jets 1--3 at B-departure
exits, jet~4 at the first distal lobe), the deflection is
$\Delta\theta_{\rm lobe}\approx16$--$19°$ below the three-jet
event plane; under any alternating-group ordering, the
deflection is at a different discrete angle.  The generic
prediction --- a discrete out-of-plane 4th jet at a specific
predicted angle rather than a broad in-plane distribution ---
is robust to the specific group ordering.  The anomaly in the
Bengtsson-Zerwas angle at $m_3+m_4>10\,{\rm GeV}$ reported by
Jeremie et al.~\cite{Jeremie2006} is consistent with this
prediction (Remark~\ref{P18:rem:jeremie}).

\item \textbf{Jet multiplicity upper bound at 12.}
The topological upper bound of 12 jets per $\QH=3$ baryon
(Proposition~\ref{P18:prop:jet_bound}) predicts that at high enough
centre-of-mass energy the jet multiplicity distribution from
identified baryons drops to zero above 12, rather than continuing
to grow as $\alpha_s^{(N-2)}$ as standard QCD would predict.
At LEP energies this is invisible (midpoint and lobe jets are
below resolution threshold), but distinguishable at future
lepton colliders above $\sim 500\,{\rm GeV}$.

\end{enumerate}

%══════════════════════════════════════════════════════════════════════════════
\section{Open Problems}
\label{P18:sec:open}
%══════════════════════════════════════════════════════════════════════════════

\begin{enumerate}[leftmargin=2em,label=\textbf{O\arabic*.}]

\item \textbf{Multi-component Jones polynomial framework.}
The topological selection rules of Section~\ref{P18:sec:reaction_rules}
require a systematic treatment of Jones polynomials for
multi-component initial and final states, including the
disjoint-union formula and the Skein-algebra action on reaction
channels. Deferred to future work.

\item \textbf{Computation of $\beta_3$ from first principles.}
Proposition~\ref{P18:prop:botner_analog} estimates $\beta_3\approx0.58$
from the A/B arc-segment energy ratio.  A derivation from the
density-feedback action would require computing the decay-plane
orientation distribution given the condensate energy profile.

\item \textbf{Quark mass from midpoint-breaking energetics.}
Conjecture~\ref{P18:conj:six_quarks} requires deriving the energy
cost $\delta J_4$ of each midpoint-site breaking.

\item \textbf{Confirmation of the Jeremie anomaly and $\Delta\theta$ measurement.}
A dedicated analysis of the BZ angle distribution in
four-jet events from baryon-initiated production, isolating the
$m_3+m_4>10\,{\rm GeV}$ region and comparing to the
$16$--$19°$ out-of-plane prediction of
Proposition~\ref{P18:prop:fourjet_oop}, would constitute
a direct test of the energy cascade of
Proposition~\ref{P18:prop:cascade_p18}.

\item \textbf{Isospin from first principles.}
Conjecture~\ref{P18:conj:isospin} connects isospin to the $\Ztwo$
writhe asymmetry; a derivation from WZW boundary conditions
is required.

\item \textbf{Quantitative production thresholds for $\QH=3$.}
Remark~\ref{P18:rem:production_conditions} identifies a two-threshold
structure (kinetic energy and collision/number-density rate)
governing where $\QH=3$ objects can form, but only qualitatively.
A quantitative treatment requires: (a) computing
$\Delta E_{2+2}=E(\QH{=}3)+E(\QH{=}1)-2E(\QH{=}2)$ and
$\Delta E_{2+1}=E(\QH{=}3)-E(\QH{=}2)-E(\QH{=}1)$ on a common
physical scale from the mass formulas of Paper~XVII~\cite{Paper17};
(b) an estimate of the cross-section $\langle\sigma v\rangle$ for
each channel from the condensate action, to convert the rate
threshold into a number-density requirement; and (c) a check that
the resulting threshold conditions correctly reproduce the QCD
deconfinement temperature $T_c\sim150$--$160\,{\rm MeV}$ as the
condition under which both thresholds are simultaneously satisfied
in the early universe.

\end{enumerate}

%══════════════════════════════════════════════════════════════════════════════
\section{Summary}
\label{P18:sec:summary}
%══════════════════════════════════════════════════════════════════════════════

This paper develops the dynamical content of $\QH=3$ confinement
through five main results.

\begin{enumerate}[noitemsep]
\item \textbf{Jones polynomial arithmetic: fusion, not sector, splits
  rational from irrational.}
The satellite Jones invariant of the trefoil
(Theorem~\ref{P18:thm:trefoil_satellite_jones}) satisfies
$|\Jones_{\rm sat}(q_5)|=1/\vp$, the \emph{same} modulus as the Hopf
link's own invariant $|\Jones_{T(2,2)}(q_5)|=1/\vp$; the two sectors
are related by the single phase $q_5^{-1}$, not by an inversion of
the golden ratio (Remark~\ref{P18:rem:richer_criterion}).  The
satellite twist family (Proposition~\ref{P18:prop:satellite_twist_family})
shows instead that rationality tracks whether the cable strands have
\emph{fused} into a genuine torus knot: every fused (odd twist count)
configuration computed is irrational, every unfused (even twist
count) configuration computed is rational.  Whether the physical
formation path passes through the rational unfused intermediate at
$n=4$, and whether this accounts for the absence of a unique
irrational BPS fixed point, is recorded as an open conjecture
(Remark~\ref{P18:rem:saddle_difficulty}), not established.
The simple-current condition $\dimq{3/2}=1$ remains the independent,
established confinement criterion of Paper~XVII; it is numerically
distinct from $|\Jones_{\rm sat}(q_5)|=1/\vp$ ($1\ne1/\vp$), and no
identification between the two is asserted here.

\item \textbf{Pion emission is topologically forced.}
$\QH=2+\QH=2\to\QH=3+\QH=1$ is the unique minimal charge-conserving
reaction forming a baryon from two $\QH=2$ objects
(Theorem~\ref{P18:thm:pion_forced}).

\item \textbf{Crossing vertices are topological lenses; energy ordering
  is 3-3-3-3.}
The $\rho_{J_4}$ field is minimum at crossing vertices, maximum
in A-segment bodies.  Energy does not exit \emph{from} crossings;
it arrives via A-segments, is redirected at crossings, and exits
at B-segment departure zones (Proposition~\ref{P18:prop:cascade_p18}).
The full jet energy ordering (hardest to softest) is
four groups of three, under the assumption that B-departure exits
dominate the primary emission channel:
(1)~primary B-departure exits (jets 1--3: hard, coplanar at $z=+r_0$);
(2)~distal lobe secondaries (jets 4--6: moderate energy,
$\approx16^\circ$--$19^\circ$ off the three-jet plane);
(3)~crossing approach-side entry disturbances (jets 7--9: softer);
(4)~midpoint torsion-twist residuals (jets 10--12: softest).
The midpoints-last portion of this ordering is a robust cascade
observation (Figure~\ref{P18:fig:gf_n256_cascade}); the specific
ordering of the middle two groups depends on which emission
mechanism dominates.  In any group-alternating ordering, the 4th
jet is topologically distinct from jets 1--3 and out of the
three-jet plane (Proposition~\ref{P18:prop:fourjet_oop}).
The topological upper bound is 12 jets per baryon
(Proposition~\ref{P18:prop:jet_bound}).

\item \textbf{Process-dependent $\phistar$ prediction in Botner form.}
By analogy with the Botner three-jet measurement~\cite{Botner1992},
the baryon-initiated distribution is
$(1+\beta_3\cos 3\phistar)/2\pi$ with
$\beta_3\approx0.58$ from the A/B arc-segment energy ratio
(Proposition~\ref{P18:prop:botner_analog}).
Photon-initiated gives $(1+\beta\cos 2\chi)/2\pi$.
These differ in angular frequency (3 vs 2) and have no
common parameters.  The Jeremie anomaly at $m_3+m_4>10\,{\rm GeV}$
in four-jet BZ angle distributions~\cite{Jeremie2006} is consistent
with a topologically-distinct 4th jet at a discrete out-of-plane
angle (Remark~\ref{P18:rem:jeremie}).

\item \textbf{Speculative correspondences.}
The $3+3$ midpoint-breaking six-quark combinatorial
(Conjecture~\ref{P18:conj:six_quarks}) and the $\Ztwo$ writhe
isospin assignment (Conjecture~\ref{P18:conj:isospin}) are
recorded as falsifiable conjectures.
\end{enumerate}

The complete sector hierarchy, including the $\QH=1$ Majorana
prediction of Paper~XVII~\cite{Paper17}, is summarised by the quantum
dimensions $\dimq{1/2}=\dimq{1}=\vp$, $\dimq{3/2}=1$
(Paper~XVII, Proposition~7.2).

%══════════════════════════════════════════════════════════════════════════════
\begin{thebibliography}{99}

\bibitem{Paper1}
F.~Manfredi,
\textit{The Density-Feedback Faddeev--Niemi Hopfion:
  Exact Algebraic Predictions for Standard-Model Parameters},
Zenodo (2026).
\href{https://doi.org/10.5281/zenodo.19342027}{doi:10.5281/zenodo.19342027}

\bibitem{Paper3}
F.~Manfredi,
\textit{Two Constants from One Knot: The Fine Structure Constant and
  the Linking Scale as Exact Predictions of the Icosahedral WZW Condensate},
Zenodo (2026).
\href{https://doi.org/10.5281/zenodo.19478629}{doi:10.5281/zenodo.19478629}

\bibitem{Paper5}
F.~Manfredi,
\textit{Colour Confinement, the QCD Scale, and Hadronic Structure
  from the Density-Feedback Hopfion},
Zenodo (2026).
\href{https://doi.org/10.5281/zenodo.19638857}{doi:10.5281/zenodo.19638857}

\bibitem{Paper15}
F.~Manfredi,
\textit{The $Q_H=3$ Sector of the Density-Feedback Hopfion},
Zenodo (2026).
\href{https://doi.org/10.5281/zenodo.20691001}{doi:10.5281/zenodo.20691001}

\bibitem{Paper16}
F.~Manfredi,
\textit{Quark Generation Masses in the Density-Feedback Hopfion:
A Survey of Ruled-Out Mechanisms and the $Q_H=3$ Gradient-Flow Landscape},
Preprint (2026).
\href{https://doi.org/10.5281/zenodo.21012650}{doi:10.5281/zenodo.21012650}

\bibitem{Paper17} F.~Manfredi,
\textit{The $Q_H=1$ Sector of the
Density-Feedback Hopfion: Topology, Group Structure, Colour Exclusion},
Zenodo (2026).
\href{https://doi.org/10.5281/zenodo.21013412}{doi:10.5281/zenodo.21013412}

\bibitem{Jones1985}
V.~F.~R.~Jones,
\textit{A polynomial invariant for knots via von Neumann algebras},
Bull.\ Amer.\ Math.\ Soc.\ \textbf{12}, 103--111 (1985).
\href{https://doi.org/10.1090/S0273-0979-1985-15304-2}{doi:10.1090/S0273-0979-1985-15304-2}

\bibitem{Kauffman1987}
L.~H.~Kauffman,
\textit{State models and the Jones polynomial},
Topology \textbf{26}, 395--407 (1987).
\href{https://doi.org/10.1016/0040-9383(87)90009-7}{doi:10.1016/0040-9383(87)90009-7}

\bibitem{Witten1989}
E.~Witten,
\textit{Quantum field theory and the Jones polynomial},
Commun.\ Math.\ Phys.\ \textbf{121}, 351--399 (1989).
\href{https://doi.org/10.1007/BF01217730}{doi:10.1007/BF01217730}

\bibitem{Rolfsen1976}
D.~Rolfsen,
\textit{Knots and Links},
Publish or Perish, Berkeley (1976).
ISBN 978-0914098164

\bibitem{BirmanBrendle2005}
J.~S.~Birman and T.~E.~Brendle,
\textit{Braids: A Survey},
in \textit{Handbook of Knot Theory} (Elsevier, 2005), 19--103.
\href{https://doi.org/10.48550/arXiv.math/0409205}{doi:10.48550/arXiv.math/0409205}

\bibitem{Faddeev1997}
L.~D. Faddeev and A.~J. Niemi,
\textit{Stable knot-like structures in classical field theory},
Nature \textbf{387}, 58 (1997).
\href{https://doi.org/10.1038/387058a0}{doi:10.1038/387058a0}

\bibitem{Botner1992}
O.~Botner et al.\ (DELPHI Collaboration),
\textit{Study of orientation of three-jet events in $Z^0$ hadronic
decays using the DELPHI detector},
Phys.\ Lett.\ B \textbf{284}, 177--187 (1992).
\href{https://inspirehep.net/literature/338132}{https://inspirehep.net/literature/338132}

\bibitem{Jeremie2006}
H.~Jeremie, P.~Leblanc, E.~Lefebvre, D.~Davignon and G.~Karapetian,
\textit{Remarks concerning the study of four-jet events from hadronic
decays of the $Z^0$},
arXiv:hep-ex/0611048 (2006).
\href{https://doi.org/10.48550/arXiv.hep-ex/0611048}{doi:10.48550/arXiv.hep-ex/0611048}

\bibitem{Kleckner2013}
D.~Kleckner and W.~T.~M.~Irvine,
\textit{Creation and dynamics of knotted vortices},
Nature Physics \textbf{9}, 253--258 (2013).
\href{https://doi.org/10.1038/nphys2560}{doi:10.1038/nphys2560}

\bibitem{Amaldi2005}
J.~Abdallah et al.\ (DELPHI Collaboration),
\textit{Measurement of the energy dependence of hadronic jet rates and
the strong coupling $\alpha_s$ from the four-jet rate with the DELPHI
detector at LEP},
Eur.\ Phys.\ J.\ C \textbf{38}, 413--426 (2005).
\href{https://doi.org/10.1140/epjc/s2004-02060-7}{doi:10.1140/epjc/s2004-02060-7}

\bibitem{Abe1997}
K.~Abe et al.\ (SLD Collaboration),
\textit{Measurement of leading particle effects in decays of $Z^0$
bosons into light flavors},
Phys.\ Rev.\ Lett.\ \textbf{78}, 3442--3446 (1997).
\href{https://doi.org/10.1103/PhysRevLett.78.3442}{10.1103/PhysRevLett.78.3442}


\bibitem{Acton1993}
P.~D.~Acton et al.\ (OPAL Collaboration),
\textit{A determination of $\alpha_s(M_{Z^0})$ at LEP using resummed
QCD calculations},
Z.\ Phys.\ C \textbf{59}, 1--20 (1993).
\href{https://doi.org/10.1007/BF01555834}{10.1007/BF01555834}

\bibitem{PDG2024}
S.~Navas et al.\ (Particle Data Group),
\textit{Review of Particle Physics},
Phys.\ Rev.\ D \textbf{110}, 030001 (2024).
\href{https://doi.org/10.1103/PhysRevD.110.030001}{doi:10.1103/PhysRevD.110.030001}

\bibitem{McKay1980}
J.~McKay,
\textit{Graphs, singularities, and finite groups},
Proc.\ Symp.\ Pure Math.\ \textbf{37}, 183--186 (1980).
\href{https://doi.org/10.1090/pspum/037}{doi:10.1090/pspum/037}

\bibitem{Verlinde1988}
E.~Verlinde,
\textit{Fusion rules and modular transformations in 2D conformal field theory},
Nucl.\ Phys.\ B \textbf{300}, 360--376 (1988).
\href{https://doi.org/10.1016/0550-3213(88)90603-7}{doi:10.1016/0550-3213(88)90603-7}

\end{thebibliography}

\end{document}